\section{Finite tests and corrections to the manuscript}\label{sec:F}
Throughout this section, $\Cat_{\mathrm{hot}}$ means functors modulo finite
zigzags of natural transformations, and $\mathrm{Hot}$ means the ordinary
unpointed homotopy category of spaces. These are different categories.
The supplied proofs are referred to by their original identifiers
\textsc{INDLOCAL} B.8, B.9, C.22, C.25, C.26, C.28. The following checks
retain their specified comparison maps.

\subsection{Finite stages: the checks of B.8, B.9, and the colimit assertion of C.22}
\begin{proposition}[B.8 at the finite bound]\label{F:B8}
Let $\mathcal C$ be a skeleton of finite categories in
$\Cat_{\mathrm{hot}}$, and localize it at nerve equivalences. A presheaf
$F$ on this localization is the restriction of an ind-object invariant
on finite tests if and only if its pullback has the following properties:
$F(\varnothing)=*$, it takes finite disjoint unions to products, and,
for every $u,v:A\rightrightarrows B$, the map
\[
 F(C)\longrightarrow\operatorname{Eq}(F(B)\rightrightarrows F(A))
\]
is onto, where $C$ is the finite Grothendieck homotopy coequalizer of
$u,v$. Every such restriction admits finite stages. This does not assert
global invariance of that presentation on arbitrary small tests.
\end{proposition}
\begin{proof}
\textbf{Check of necessity, first step.} A representable takes an empty
source to one element and a coproduct of sources to a product. Filtered
colimits of sets commute with finite products: representatives for a
finite tuple and a finite list of equality witnesses occur together at
one stage. This proves the first two conditions, including the empty case.

\textbf{Check of necessity, the homotopy coequalizer step.} Represent an
equalizing element by $b:B\to K_j$ at one stage, increasing $j$ until
$bu,bv$ are joined by a finite transformation zigzag. Let $Z$ be its
oriented interval. Use the free finite acyclic indexing category on $Z$
and a further vertex $b$, with one arrow from each endpoint to $b$.
Place $A$ at the vertices of $Z$, identity functors on its edges, $B$ at
$b$, and $u,v$ on the endpoint arrows. Its Grothendieck construction
$C_Z$ is finite. The actual zigzag and $b$ define a functor $C_Z\to K_j$.
Collapsing $Z$ gives $C_Z\to C$, equal to the identity on the copy of
$B$. On the double mapping cylinders this collapses the subdivided
interval between the two attaching maps; it is a homotopy equivalence
relative to $BB$. Thus finite-test invariance transports the element to
$F(C)$ with exactly the required restriction to $F(B)$. This verifies the
comparison direction in the supplied proof; no homotopy between functors
is silently replaced by one transformation.

\textbf{Check of sufficiency.} The category of elements of $F$ is nonempty
by the empty-category value. The coproduct condition gives common upper
bounds for two elements. For two parallel arrows represented by $u,v$,
the last condition gives an element over $C$ extending the common element
over $B$. The two tagged identity arrows in the Grothendieck construction
give transformations from the copy of $A$ to $wu$ and $wv$; consequently
$wu=wv$ in $\Cat_{\mathrm{hot}}$. Thus the element category is filtered.
Its objects form a set, since the finite-category skeleton and every
presheaf value are sets. The usual element-by-element Yoneda map
$\colim_{(A,a)}\Hom(-,A)\to F$ is onto by its universal element and
injective by the defining element-category arrows. This is the required
finite-stage presentation and completes all steps of B.8 used here.
\end{proof}

\begin{proposition}[B.9, including its strict idempotent step]\label{F:B9}
A filtered ind-object of $\Cat_{\mathrm{hot}}$ with at most $m$ arrows
at every stage is isomorphic to a constant category with at most $m$
arrows.
\end{proposition}
\begin{proof}
\textbf{Step 1: finite reduction.} There are finitely many category tables
on at most $m$ arrows, and finitely many functors between any two.
Choose one representative of each isomorphism class and form the finite
full subcategory $\mathcal D$ in $\Cat_{\mathrm{hot}}$. Transport every
stage into $\mathcal D$. The case $m=0$ is the constant empty category.

\textbf{Step 2: finite element category.} For $A\in\mathcal D$, the value
$F(A)=\colim_i\Hom(A,K_i)$ has cardinality at most
$\max_{T\in\mathcal D}|\Hom(A,T)|$. Otherwise finitely many distinct
colimit elements would already violate that bound at a common stage.
The element category of $F|_{\mathcal D}$ is therefore finite and is
filtered by the same common-stage and equality-witness argument.

\textbf{Step 3: the actual retract.} A cocone on its finite identity
diagram has a vertex $(T,x)$ and arrows $t_d:d\to(T,x)$ satisfying
$t_e t_d=t_d$, including $t_e^2=t_e$. Applying the density presentation
gives $h_T\to F\to h_T$ whose first composite is $1_F$ and whose second
is $h_{t_e}$. The full inclusion of $\mathcal D$ induces a fully faithful
inclusion of ind-categories: both Hom sets are the same limit of filtered
colimits of the same Hom sets. This is a retract of the original
ind-object, not just an isomorphism on the finitely many chosen tests.

\textbf{Step 4: splitting inside finite categories.} Choose an actual
endofunctor $f:T\to T$ for $t_e$. Powers of $f$ are eventually periodic.
Choose $r$ beyond the preperiod and divisible by the period. Then
$f^{2r}=f^r$ strictly, while $[f^r]=[f]$ since $[f]$ is idempotent.
The objects and arrows fixed by $f^r$ form a category $L$, with inclusion
$j:L\to T$ and $s:T\to L$ satisfying $sj=1$ and $js=f^r$.
It has at most $m$ arrows. The two splittings of the same idempotent
identify $F$ with $h_L$ in the full ind-category. This confirms the
supplied proof without confusing a homotopy idempotent with a strict one.
\end{proof}

\begin{proposition}[C.22: exactly the strict colimit assertion]\label{F:C22}
For full subcategory inclusions $K_0\subset K_1\subset\cdots$ and a
category $X$ with finitely many objects, the canonical functor
\[
 \colim_n\operatorname{Fun}(X,K_n)\longrightarrow
 \operatorname{Fun}(X,\colim_n K_n)
\]
is an isomorphism. For an arbitrary filtered diagram, it is an
isomorphism if $X$ has finitely many arrows. The corresponding assertion
on components holds. These statements do not identify components with
derived mapping classes.
\end{proposition}
\begin{proof}
The images of the finitely many objects lie at a common stage. Fullness
puts every arrow image there, and faithfulness preserves every equation,
even if $X$ has infinitely many arrows. A natural transformation has only
finitely many object-components; fullness and faithfulness give the same
argument for its naturality equations. For a general filtered diagram,
use instead the finite list of arrow images and the finite identity and
composition relations of $X$, and move to a common stage where they all
hold. Equality is handled by the same finite list. Finally every zigzag
witnessing a component equality is finite. These are respectively the
surjectivity, injectivity, and composition checks in C.22.
\end{proof}

\subsection{The finite obstruction of C.25 and the finite example of C.26}
\begin{proposition}[C.25 checked by explicit prisms]\label{F:C25}
Weak-equivalence invariance of the constant-endpoint localization functor
can fail for a finite relative category, even with all arrows marked,
and for an equivalence between finite test categories.
\end{proposition}
\begin{proof}
\textbf{The source equivalence.} Let $Y$ be the four-point crown
$r_0,r_1<q_0,q_1$, let $M$ be the six-point poset obtained by adjoining
two maxima $s,n$ above all of $Y$, and mark every arrow of $M$.
Let $B$ be the parallel-pair category $a\rightrightarrows b$, with arrows
$f,g$. Send the minima of $Y$ to $a$, its maxima to $b$, the edge
$r_0q_1$ to $g$, and the other edges to $f$. This functor $u:Y\to B$
sends the generating cycle
\[
 \gamma=[r_0q_0]-[r_1q_0]+[r_1q_1]-[r_0q_1]
\]
to $f-g$. Both realizations are circles, so $Bu$ is a homotopy equivalence.

\textbf{The vanishing source value.} Because $M$ is a poset,
$M^B\cong M^{[1]}$. Source evaluation has the constant-arrow section,
and the arrow itself is a natural transformation from that section
followed by evaluation to the identity. Hence $B(M^B)\simeq BM=S^2$.
Marking all arrows makes ordinary localization the fundamental groupoid,
so $H_{s,s}(B)=\pi_1B(M^B,s_B)=0$.

\textbf{The nonzero target value.} With $j:Y\hookrightarrow M$, the loop
\[
 s_Y\longleftarrow j\longrightarrow n_Y
       \longleftarrow(r_0)_Y\longrightarrow s_Y
\]
lies in $B(M^Y)$. Evaluation on $BY$ sends its product with $\gamma$
to $n*\gamma-s*\gamma$. On each of the four edge prisms the nondegenerate
triangle is the indicated cone triangle; the constant return legs give
zero. The resulting eight triangles have coefficients $\pm1$, boundary
zero, and form the primitive fundamental cycle of the octahedral sphere
$BM$. Evaluation on that cycle is additive under concatenation, giving
a homomorphism $\pi_1B(M^Y,s_Y)\to\Z$ taking the loop to one.
Thus restriction from $H_{s,s}(B)$ is not onto. The independent program
\texttt{f\_indlocal\_checks.py} checks the entire relation table, all four
transformations, all prism coefficients and boundaries, and the counts
$198$ and $18$ of functors into $M$. This confirms each finite step of
C.25, rather than accepting the cited PATHS result as an authority.
\end{proof}

\begin{proposition}[The five-arrow counterexample in C.26]\label{F:C26-finite}
There is a nontrivial finite pointed relative category satisfying all
three printed tests at IX ms.~22--23, including global invariance and
both loop and suspension formulas.
\end{proposition}
\begin{proof}
Take the full category $E$ on the pointed sets of cardinalities one and
two, and mark only identities. Its arrows are
\[
 1_0,\quad 1_x,\quad u:0\to x,\quad v:x\to0,\quad z=uv:x\to x,
 \qquad vu=1_0,\quad z^2=z\ne1_x.
\]
These actual set maps specify every composite. The program verifies all
$34$ composable triples. The object $0$ is both initial and terminal.

Localization at identities changes neither $E$ nor any $E^X$.
A transformation between two constant diagrams chooses one arrow on
each connected component, and the naturality equation says precisely
that these choices are constant along every zigzag. Therefore
\[
 H_{a,b}(X)=\Hom_E(a,b)^{\pi_0X}.
\]
This includes the empty category and proves global invariance. The
localized mapping spaces are the discrete sets $\Hom_E(a,b)$.
All their based loop spaces are contractible; their trivial fundamental
groups are represented in both variables by $0$. Set
$\Sigma a=0=\Omega b$. Both mapping-space formulas then have contractible
sides, with their unique group coherences. In the explicit mixed path
models, every backward edge is an identity. Deleting those edges and
repeating their vertices are inverse functors; the naturality square at
an identity forces equality of the repeated components. These
isomorphisms respect padding. To check the categorical contraction
after padding, regard a positive path as an eventually constant chain,
equal to $a$ on its left tail and $b$ on its right tail. Fix one cut.
The functor $T$ keeps the chain to the left of the cut, replaces the
right tail by $b$, and uses the remaining composite at the cut.
Tail composites give $1\Rightarrow T$. If $S$ inserts the total effect
at the cut, prefix composites give $S\,\mathrm{effect}\Rightarrow T$.
These transformations are uniform on the union and preserve the
endpoints. For one-sided padding the same construction uses the fixed
endpoint cut. Hence effect and $S$ are categorical homotopy inverses
for the actual positive, mixed and symmetric padding systems in this
example. Thus the stronger path-model reading also gives the discrete
Hom sets. Finally $1_x\ne z$ survives
localization, so the resulting category is not equivalent to a point.
\end{proof}

\subsection{V ms.~51--52: the class of finite-category types and its tests}
\begin{proposition}\label{F:Hot0}
Let $\mathrm{Hot}_0$ be the full subcategory of homotopy types represented
by finite categories. It contains every finite CW homotopy type and
strictly more, but does not contain $K(\Z,2)$, $K(\Z/2,2)$, or $\Omega S^3$.
For connected finite CW complexes, its one-object part consists exactly
of the simply connected types. Restricted unpointed set-valued tests
$X\mapsto([BC,X])_{C\text{ finite}}$ are neither faithful nor conservative.
\end{proposition}
\begin{proof}
A finite CW complex has the homotopy type of a finite simplicial complex;
its face poset is a finite category with that type. The two-arrow group
$C_2$ has classifying space $\RP^\infty$, outside finite CW type. The
exclusions follow from Corollaries~\ref{A:rational-nilpotence} and
\ref{A:K22}. The one-object assertion is Theorem~\ref{A:bounded} and the
checked constructive theorem in Section~\ref{sec:B}.

For the test assertion, let $G$ be the group of finitary permutations of
$\N$ and let $s:G\to G$ insert a fixed zero:
$s(\sigma)(0)=0$, $s(\sigma)(i+1)=\sigma(i)+1$. It is injective but misses
$(0\ 1)$, so $Bs$ is not a homotopy equivalence. For any finitely generated
$P$ and $\rho:P\to G$, the images of generators have support in some
$\{0,\ldots,m\}$. The finite cycle $c=(0\ 1\ \cdots\ m+1)$ satisfies
$c\rho(p)c^{-1}=s\rho(p)$ for every $p$. The identity is verified by
following each point inside and outside this finite support.

For connected $X$, maps to $BG$ are homomorphisms $\pi_1X\to G$ modulo
conjugacy: construct on edges, use relations on triangles, and extend
uniquely through higher cells because $BG$ is a $1$-type. Each finite
category has finitely many components and finitely presented fundamental
groups, by its finite arrow presentation. Thus $Bs$ induces exactly the
identity on every $[BC,BG]$. Nevertheless $Bs\ne1$ in $\mathrm{Hot}$,
since a homomorphism conjugate to the identity is surjective. This is a
single counterexample to both faithfulness and conservativity. The
finite permutation check in \texttt{f\_manuscript\_models.py} is independent
of the supplied program; the support argument proves the unbounded claim.
\end{proof}

\subsection{VII ms.~121 and 128: compactness and finite calculations}
\begin{proposition}\label{F:noncompact}
Finiteness of the arrow set does not make $BC$ a compact homotopy test.
In particular the canonical comparison
\[
 \colim_r[\RP^\infty,\RP^r]\longrightarrow
 [\RP^\infty,\RP^\infty]
\]
is not surjective, although $\RP^\infty=BC_2$ for a finite category.
\end{proposition}
\begin{proof}
The homotopy colimit of the skeletal inclusions is $\RP^\infty$.
If its identity factored through $\RP^r$ in the homotopy category,
the identity on $H^{r+1}(-;\F_2)$ would factor through zero. But
$H^*(\RP^\infty;\F_2)=\F_2[t]$ with $|t|=1$. This contradiction gives
the stated failure, not merely infinite dimensionality of one nerve.
Finite compatible triangulations of the successive real projective
spaces give the same example using finite face-poset stages and the
specified maps of realizations. Proposition~\ref{F:C22} concerns strict
functor categories and therefore does not contradict this example.
\end{proof}

\begin{proposition}\label{F:computation}
For a finite category $C$ and a fixed degree $n$, integral homology
$H_n(BC;\Z)$ is determined by a finite exact integer calculation.
The fundamental group has an effectively given finite presentation.
Neither statement is an algorithm deciding all homotopy questions about
finite categories, in particular triviality of their fundamental groups.
\end{proposition}
\begin{proof}
The normalized nerve has finitely many nondegenerate simplices in each
fixed degree. Its boundary is the finite alternating sum of deletion and
composition faces. Smith normal forms of the matrices through degree
$n+1$ determine $H_n$, with no inference about larger degrees. Choose a
spanning forest of the finite underlying graph. The remaining arrows
generate $\pi_1$, and the finite composition triangles give the relations.
Every finitely presented group occurs as $\pi_1$ of a finite presentation
complex, hence of a finite face-poset model. An algorithm deciding
triviality for all finite-category fundamental groups would therefore
decide triviality for every finite group presentation. Adian's
1957 Theorem~1, read in the primary English translation
\cite[p.~73]{adian1957translation}, states that for every fixed finitely
presented group $F_0$ its isomorphism-recognition problem on finite
presentations is undecidable. Set $F_0=1$; it is finitely presented, so
all hypotheses hold and the proposed algorithm is impossible. The
conversion to a finite simplicial presentation complex and then its face
poset is effective, so this is a reduction of algorithms. No impossibility
is asserted for the usual restricted algorithms for simply connected
complexes.
\end{proof}

\subsection{XII Appendix ms.~8--11 and 15: the actual small categories}
\begin{proposition}\label{F:XII-models}
The parallel-pair category has $4$ arrows and type $S^1$. The category
$Z$ drawn at XII Appendix ms.~10--11 has $4$ objects, $10$ arrows, and
type $S^2$, with $\pi_1(BZ)=0$ and $\pi_2(BZ)=\Z$. It does not retract
onto its parallel-pair subcategory. The least arrow counts for
$S^1,S^2,S^3$ are $4,5,7$, respectively; the $S^3$ models of size seven
are exactly the four monoids in Section~\ref{sec:C}.
\end{proposition}
\begin{proof}
The nerve of the parallel pair is a graph with two vertices and two
edges, hence a circle. For $Z$, name the objects $0,a,b,c$. There are
two arrows $0\to a$, one each $a\to b,a\to c$, and a single composite in each of $\Hom(0,b)$ and $\Hom(0,c)$:
the two paths to each of these targets are identified. There are no other
nonidentities. This gives $4+2+1+1+1+1=10$ arrows. The full subcategories
on $0,a,b$ and $0,a,c$ are cones on the parallel pair (their terminal
objects are $b,c$). Their nerves intersect in that pair and their union
is all of $NZ$. Thus $BZ$ is its unreduced suspension, homeomorphic to
$S^2$. A retraction onto the circle would inject its $H_1$ into zero,
which is impossible. This confirms the manuscript's own corrective
footnote 7 and identifies the first false step of the earlier proposed
retraction argument. All compositions and the integral nerve homology
are checked in \texttt{f\_manuscript\_models.py}. The minimum counts and
classification follow from the independent enumeration and resolutions
of Section~\ref{sec:C}, not from this larger model.
\end{proof}

\begin{corollary}[The replacement at XII Appendix ms.~15]\label{F:XII15}
For the displayed category $V'_0=Z$, replace the claim
``$\pi_1\cong\Z$'' by $\pi_1=0$, $\pi_2\cong\Z$.
The proposed reason for non-asphericity for $W\subseteq W_1$ is false.
A valid sufficient containment for this particular example is
$W\subseteq W_2$ (in particular $W\subseteq W_\infty$).
\end{corollary}
\begin{proof}
The preceding calculation applies to the very same diagram.
For $W=W_1$, the map $Z\to *$ belongs to $W_1$: it is bijective on
components and an isomorphism on fundamental groups. Thus the old
parenthetical conclusion actually fails for an allowed choice of $W$.
For $W\subseteq W_2$, membership of $Z\to *$ would force its nonzero
second homotopy group to vanish, so it is not $W$-aspherical. This repairs
the implication at the level the computed homotopy type supports.
\end{proof}

\subsection{Checking the type prerequisites of C.20 and C.22}
The supplied \textsc{INDLOCAL} extract starts at C.20. Its references to
C.1--C.19 are therefore not proofs present in that extract. We supply
the precise prerequisites used below, and then check the written first
proof of C.20. This avoids treating an unavailable earlier argument as
an authority.

\begin{lemma}[The type-map contraction]\label{F:type-contraction}
A type is a finite word $\sigma$ in $+,-$, including the empty word.
Its interval $I_\sigma$ is the free category on the oriented line with
vertices $0,\ldots,|\sigma|$. Maps of types are endpoint-preserving,
nondecreasing functors of these intervals. For fixed $\sigma,\tau$, the
category $P(\sigma,\tau)$ of such functors and endpoint-fixed natural
transformations is either empty or categorically contractible: its
identity is connected to a constant functor by a finite zigzag of
natural transformations.
\end{lemma}
\begin{proof}
Write $n=|\sigma|$, $m=|\tau|$. A type map is equivalent to a weakly
increasing allocation
\[
 a_1\leq\cdots\leq a_m,\qquad
 a_j\in\{1,\ldots,n\},\quad \sigma_{a_j}=\tau_j.
\]
Target edge $j$ is traversed by source edge $a_j$. Conversely the image
of source vertex $i$ is the number of allocated target edges with
$a_j\leq i$. Each source edge then traverses a string of its own sign
or an identity, so this is exactly the required functor. This also
handles an empty target; a nonempty target cannot be allocated to an
empty source.

The interval category is a poset. A transformation from allocation
$a$ to allocation $b$ therefore exists uniquely precisely when
\[
 a_j\geq b_j\quad(\tau_j=+),\qquad
 a_j\leq b_j\quad(\tau_j=-).
\]
Indeed a component crosses edge $j$ in the positive vertex direction
exactly when its allocation has moved earlier, and can cross it in that
direction exactly when its sign is positive. The reverse statement
holds for negative edges. Pointwise existence implies naturality in a
poset, so these inequalities give both necessity and sufficiency.

Suppose allocations exist. Let $q_j$ be the minimum of their $j$th
coordinates. Taking these minima preserves monotonicity and the required
colors, so $q$ is an allocation.
Let $F_j(a)$ replace the first $j$ coordinates of $a$ by those of $q$.
The junction remains increasing, since $q_j\leq a_j\leq a_{j+1}$.
The displayed inequalities show that every $F_j$ is an order-preserving
endofunctor. For a positive $\tau_j$ there is
$F_{j-1}\Rightarrow F_j$; for a negative $\tau_j$ there is the reverse
transformation. Since $F_0=1$ and $F_m$ is constant at $q$, this is the
claimed contraction. No connectedness or coherence assertion is left
to a choice of individual paths.
\end{proof}

\begin{proposition}[The uniform fibred comparison used in C.28]
\label{F:C23-uniform}
For $u:J\to T$, the map of fibred Grothendieck constructions
\[
 \int_J^{\mathrm{fib}}F u^{\op}\longrightarrow
 \int_T^{\mathrm{fib}}F
\]
is a nerve weak equivalence for every strict $F:T^{\op}\to\Cat$
if and only if each comma $u/t$ is weakly contractible. Parallel
pullback homotopies and a cofinal reachability sequence alone do not
imply this condition. In particular they do not justify the universal
comparison proposed at VII ms.~99--110.
\end{proposition}
\begin{proof}
For $(t,z)$ in the target total category, forget fibre data in its
comma category. The resulting functor to $u/t$ has the section sending
$(j,\alpha:u(j)\to t)$ to the identity of $F(\alpha)z$.
The fibre data of any object are exactly an arrow
$x\to F(\alpha)z$. This arrow gives a natural transformation from
the identity to section after projection. A comma arrow $v$ obeys
$\alpha=\alpha'u(v)$, and contravariance gives the exact naturality
equation $F(\alpha)=F(u(v))F(\alpha')$. The comma spaces are therefore
equivalent to $B(u/t)$. Quillen's Theorem A proves sufficiency.
For necessity use the discrete presheaf $F(s)=\Hom_T(s,t)$: its
total category is $T/t$, with terminal object, and its pullback total
category is exactly $u/t$.

For the printed failure take $T=\mathbf BC_2$ and let
$u:\N^{\op}\to T$ send every arrow to the identity. Use the constant
terminal fibre. Parallel pullbacks agree, the reachability preorder
is a point, and a countable cofinal sequence exists. The source total
category is $\N^{\op}$ and is contractible, whereas the target is
$\mathbf BC_2$ with fundamental group $C_2$. Thus the specified map
is not a weak equivalence. For the separate smooth/proper suggestion,
the types $+-+$ and $+$ each have only identity endomorphisms, and
there are exactly two maps between them, choosing one of the two
positive edges. The corresponding incoming and outgoing lift
categories over their reachability arrow are discrete two-point
categories. They are not even connected. Lemma~\ref{F:type-contraction}
instead supplies transformations between type maps; it must not be
confused with either of these discrete lift categories.
\end{proof}

\begin{proposition}[C.22, the remaining type and colimit assertions]
\label{F:C22-types}
Concatenation of types and their interval functors is strictly
associative and unital. The empty type is terminal. The opposite type
category is neither filtered nor totally $W_\infty$-aspherical.
Nevertheless every actual projective sequence of types gives a
covariant path diagram indexed by the filtered category $\N$; its
ordinary categorical colimit has the homotopy type of its homotopy
colimit. Together with Proposition~\ref{F:C22}, these statements check
all assertions of \textup{\textsc{INDLOCAL}} C.22.
\end{proposition}
\begin{proof}
\textbf{Intervals and concatenation.} Gluing two free oriented lines at
their end and initial vertices gives their concatenated free line.
Functors out, and transformations between such functors, are exactly
pairs agreeing at the junction. Thus this is the asserted categorical
pushout and gives the corresponding pullback description of paths.
Blockwise concatenation of vertex maps is strictly associative, has the
empty block as strict unit, and commutes with composition. A type map
is uniquely a concatenation, over its source edges, of maps
$\varepsilon\to\varepsilon^l$, $l\geq0$. Under composition each such
string is refined by the lengths of the second map's strings.
Consequently the descriptions by functions $\{1,\ldots,l\}\to\N$
and their iterates in every nerve degree at VII ms.~50'--58 follow,
including zero lengths and identity insertions. Endpoint preservation
is exactly what preserves endpoint evaluations in an arrow
$(u,\eta):c\to d\circ u$ of the fibred path category.

\textbf{The restriction on interval maps.} The functor to categories
is faithful but not full onto all endpoint-preserving interval
functors. The vertex sequence $(0,1,0,1,2,3)$ defines a functor
$I_{+-+-+}\to I_{+-+}$, since its five arrow images are respectively
$0\to1,0\to1,0\to1,2\to1,2\to3$. It preserves endpoints and is not
nondecreasing. This checks the precise restriction in the written proof.

\textbf{Failure of filteredness.} By the allocation description,
$J=\mathrm{Ty}^{\op}$ is the category of finite linearly ordered
sets colored by $+,-$ and monotone color-preserving maps. The two maps
$+\rightrightarrows +-+$ selecting its first and last positions have
no common equalizing postcomposition. Such a postcomposition would
identify the two positive positions and hence also the intervening
negative position by monotonicity, contradicting color preservation.

\textbf{An explicit witness against total asphericity.} Put $p=++$
and $S(t)=J(p,t)$. The category $E=\operatorname{El}_J(S)=p\downarrow J$
has initial object $1_p$. An object is a pair of positive positions
$i\leq j$ in a word $t$. Let $U$ be the cosieve generated by the
collapse $p\to+$, and $V$ the cosieve generated by the map
$p\to+-+$ selecting the outer positions. They have the following
exact descriptions: $U$ consists of equal pairs $i=j$; $V$ consists
of pairs with a negative position strictly between them. They are
nonempty, connected, disjoint cosieves. For example an image of a
negative position strictly between two positive ones must remain
strictly between them, because colors prevent equality.

On $E$ use the two-element action outside $U\cup V$ and the singleton
action on $U\cup V$, with identity or collapse structure maps. Inducing
this action to $J$ gives the explicit covariant set functor $L$: a pair
$i\leq j$ has one copy if $i=j$ or a negative position lies between
them, and two copies otherwise. A monotone colored map sends the pair
to its image and retains its copy label unless the target has one
copy. The cosieve property proves functoriality. Its element category
is two copies of $E$ glued on $U\sqcup V$. A chain cannot leave a
cosieve after entering it, so this is also a pushout of nerves.
Both copies of $BE$ are contractible; $BU$ and $BV$ are connected.
Mayer--Vietoris consequently gives
\[
 \widetilde H_0(B\operatorname{El}L;\Z)=0,
 \qquad H_1(B\operatorname{El}L;\Z)=\Z.
\]
For a direct test of the definition, adjoin to $J$ a new source $*$
with $\Hom(*,t)=L(t)$, no arrows back, and its identity endomorphism.
Call the resulting category $J'$ and the inclusion $j:J\to J'$.
The comma $*\downarrow j$ is $\operatorname{El}L$ and is connected;
the comma at an old object has an initial object. Therefore the
ordinary presheaf left Kan extension satisfies $j_!1_J=1_{J'}$.
The nerve of this collage gives the homotopy cofiber of
$B\operatorname{El}L\to BJ$. Since $J$ has the initial empty word,
$BJ$ is contractible and $H_2(BJ';\Z)=\Z$. Thus the element category
of this Kan extension is not weakly contractible. This refutes total
$W_\infty$-asphericity under exactly the definition in
\textsc{PRETOT}, rather than inferring it from failure of filteredness.

\textbf{The actual sequential colimit.} Pullback reverses type maps,
so a projective sequence acts covariantly on paths with transitions
$K_{n-1}\to K_n$. Nerves commute with filtered categorical colimits:
an $r$-simplex and each equality between two such simplices are finite
lists of arrows and composition equations. The simplicial replacement
map from the filtered homotopy colimit to this colimit is a weak
equivalence. One can check this directly on homotopy groups: a
triangulated sphere, disk, and any homotopy have finitely many simplices;
their data and equations have a common stage, and a finite set of
stages and arrows has a cone at a later stage. These cones give the
required lifts and equality witnesses. Realization yields the stated
homotopy equivalence, independently of any cofinality of the sequence
inside all types.

The independent program \texttt{f\_type\_categories.py} checks all
$961$ pairs of words of length at most four, including $1471$ type
maps and $10771$ transformation tests, against the interval definition.
It also checks the complete colored-order category on words of length
at most three ($15$ objects, $209$ arrows) and the above action; its
element category has $30$ objects, $499$ arrows, $94404$ composable
triples and first mod-two Betti number one. These are checks on fixed
finite instances; the universal assertions use the arguments above.
\end{proof}

\begin{proposition}[The first proof of C.20, with its missing prerequisite supplied]
\label{F:C20-paths}
Let $(M,\Sigma)$ be a small relative category with all identities and
composites in its marking. Let $t:\N^{\op}\to\mathrm{Ty}$ be an actual
sequence such that every type receives a map from some $t_n$.
Write $K_n=F_{t_n}(a,b)$ for paths with marked negative edges, arbitrary
interior vertical arrows, and identity endpoint components. For every
small category $X$, including infinite and empty ones, effect of words
gives a natural bijection
\[
 \colim_n\pi_0\operatorname{Fun}(X,K_n)
 \cong
 \Hom_{M^X[\Sigma_X^{-1}]}(a_X,b_X).
\]
Here $\Sigma_X$ consists of the pointwise marked transformations.
The comparison respects constant paths, concatenation, and restriction
in $X$.
\end{proposition}
\begin{proof}
\textbf{The exponential-law step.} A path of type $\tau$ in $M^X$ is
exactly a functor $X\times I_\tau\to M$ with the stated endpoint and
marking conditions. Its ladders have exactly the endpoint-fixed
components in the definition. Hence
$F_\tau^{M^X}(a_X,b_X)=\operatorname{Fun}(X,F_\tau^M(a,b))$
canonically, on objects and arrows, and compatibly with type pullbacks.

\textbf{The component universal property.} Compose forward arrows and
inverses of backward marked arrows in the ordinary localization.
Type pullback does not change this word's effect. A positive ladder
square preserves the composite equation; a negative square does so
after cancelling its two invertible horizontal edges. No interior
vertical edge is inverted. Conversely components of all-type paths
form a category under concatenation. Selecting the outer vertices
compares two successive forward edges with their composite; collapse
compares an inserted identity. For $w:x\to y$ in the marking, compare
the backtrack $x\xrightarrow{w}y\xleftarrow{w}x$ to the constant
path at $x$ by the ladder with middle component $w:x\to y$ and
identity endpoint components. For the other backtrack use the ladder
to the constant path at $y$, again with middle component $w$.
Thus marked arrows become invertible. Evaluation in any category
inverting $\Sigma_X$ respects all these equations and gives its unique
factorization. This proves the localization universal property and
identifies the specified effect map.

\textbf{Surjectivity and injectivity at actual stages.} Each word has
one finite type, so pulling it back to a $t_n$ proves surjectivity.
An equality of components is a finite zigzag of total-path arrows.
Concatenate all its types, and choose a stage mapping to that word,
increasing the stage to dominate both original endpoints. Pull every
ladder back there. A type-change arrow may compare two pullbacks of
one path along different maps from this stage to the same type.
Lemma~\ref{F:type-contraction} connects these maps by transformations
uniformly in the path; composing and applying $\operatorname{Fun}(X,-)$
gives the required ladder zigzag even for infinite $X$. The same
argument at both endpoints compares the choices with the actual
transition maps of the original sequence. The equality therefore holds
at an actual common stage. Every step uses restriction, word effect,
or concatenation, proving the claimed naturalities and units.

\textbf{The exact marking hypothesis.} Pullback along
$-\to\varnothing$ requires marked identities; pullback along
$-\to--$ requires closure under composition. These conditions suffice
because a negative edge maps to a negative string or an identity.
For an arbitrary class of arrows, replacing it by its generated wide
subcategory preserves the ordinary localization by its universal
property and makes the type diagram defined. This confirms all steps
of the supplied first proof of C.20; its alternative proof via C.3
is not needed for this assertion.
\end{proof}

\subsection{The all-type and specified sequential path comparisons}
Throughout this subsection categories may be arbitrary small categories.
A marking is a subcategory containing all objects; it need not contain
all isomorphisms. Fix endpoints $a,b$ in $(M,\Sigma)$. Write
$\mathcal A(M;a,b)$ for the fibred all-type category: an object is
$p:I_\tau\to M$ with marked negative edges, and an arrow from $p$ to
$q:I_\sigma\to M$ is a type map $u:I_\tau\to I_\sigma$ and an
endpoint-fixed transformation $p\Rightarrow qu$. Write
$\mathcal W(M;a,b)$ for the same category with every component of this
transformation required to lie in $\Sigma$.

For an actual projective type sequence $t:\N^{\op}\to\mathrm{Ty}$,
write $\mathcal E_t(M;a,b)$ for its fibred total category. Thus an
arrow from stage $n$ to stage $m$ requires $n\geq m$ and is a
transformation $p\Rightarrow q\,t(n\to m)$. These conventions specify
the functors, not just the spaces up to an unspecified equivalence.

\begin{lemma}[The marked comparison and its marking hypothesis]
\label{F:marked-hammocks}
There is a natural weak equivalence, expressed by the canonical
Grothendieck and reduction zigzag,
\[
 B\mathcal W(M;a,b)\simeq
 \operatorname{Map}_{L(M,\Sigma)}(a,b).
\]
It respects concatenation. This statement holds for the marking
convention just specified.
\end{lemma}
\begin{proof}
The fixed-type nerves are the fixed-endpoint marked hammock diagrams.
The type category has its Reedy factorization into a vertex-surjection
followed by a vertex-injection. In the opposite index, a latching map
is the inclusion of diagrams with the relevant collapsed columns equal
to identities; intersections correspond to the generated common
quotient. Its latching maps are consequently monomorphisms. The
matching indexing categories for constant diagrams are empty or
connected: the reduced interval consisting of the endpoints and the
direction-change vertices is initial among proper injective
subintervals whenever there is one. Hence ordinary colimit computes
homotopy colimit. These are the precise combinatorial checks in
\cite[Proposition~4.8, Lemma~4.16]{arakawacnossen2026}; they use
identities and closure under composition, not inclusion of every
isomorphism.

The colimit is the reduced hammock simplicial set: identity columns
are deleted and adjacent equally directed columns are composed.
\cite[Section~2.1 and Proposition~2.2]{dwyerkan1980calculating}
identify this with the simplicial localization. The convention is
explicit in \cite[Section~3.1]{dwyerkan1980function}: the subcategory
contains all objects. Thus no enlargement of $\Sigma$ has occurred.
For the infinity-categorical universal property, one may now enlarge
to the closure $\overline\Sigma$ of \cite[Section~3.4 and
Lemma~5.1]{dwyerkan1980simplicial}; that lemma makes the induced
simplicial-localization map a weak equivalence. This closure contains
every isomorphism, so \cite[Theorem~4.7 and Corollary~1.11]{arakawacnossen2026}
applies with its stated convention. A functor to any infinity-category
which inverts $\Sigma$ also inverts $\overline\Sigma$: pass to its
homotopy category and use the defining ordinary localization.
Consequently the two infinity-localization universal properties agree.
This explains, rather than assumes, the passage between conventions.

For orientation, if $F:T^{\op}\to\Cat$, then
\[
 \left(\int_T^{\mathrm{fib}}F\right)^{\op}
   =\int_{T^{\op}}^{\mathrm{opf}}F^{\op}.
\]
Reversing a nerve simplex gives a natural homeomorphism
$B(C^{\op})\cong BC$. Apply the verified Grothendieck homotopy-colimit
theorem to the right side and these homeomorphisms to the fibres.
This is the required comparison for the fibred convention used here.
All reduction operations and these natural comparisons commute with
block concatenation, which gives the last assertion.
\end{proof}

\begin{lemma}[A relative interval contracts all arbitrary-vertical paths]
\label{F:relative-interval}
Give $I_\tau$ the marking generated by its negative edges. Every
nonempty category $F_\rho(0,|\tau|)$ of relative paths in this interval,
with arbitrary endpoint-fixed transformations, is categorically
contractible. It is nonempty exactly when there is a type map
$\rho\to\tau$. Consequently $B\mathcal A(I_\tau)$ is contractible.
If $t$ is universally cofinal in the reachability sense of C.20, then
$B\mathcal E_t(I_\tau)$ is contractible as well.
\end{lemma}
\begin{proof}
For such a path $h$, regarded as a sequence of target vertices, set
\[
 \overline h(k)=\max_{j\leq k}h(j).
\]
A path can cross a positive target edge backwards neither on a
positive source edge (there is no such target arrow) nor on a negative
source edge (its image would be unmarked). Every numerical decrease
therefore stays within a negative run. The segment
$\overline h(k)\to h(k)$ is a marked arrow. Moreover
$\overline h$ is nondecreasing and preserves the required arrow
directions: when its value increases, its new segment is a subsegment
of the arrow traversed by the corresponding edge of $h$. In particular
a negative source edge still has marked image. Endpoints are fixed.
Thus $\overline h$ is a type map and there is an endpoint-fixed
transformation $\overline h\Rightarrow h$.

Numerical maximum is order-preserving separately in each variable on
any oriented-line poset: clipping a directed segment at a fixed
vertex gives a subsegment or an identity. Hence $h\mapsto\overline h$
is a functor on the path category, and the preceding transformation
is natural in $h$. It retracts onto $P(\rho,\tau)$ and is the identity
there. Lemma~\ref{F:type-contraction} proves categorical
contractibility and the nonemptiness assertion.

The supported full subcategory of types is contractible. The
endofunctor $\rho\mapsto\rho\tau$ has natural type maps to both
$\rho$ and the constant type $\tau$, obtained by collapsing the
other block. Its values are supported, and $\tau$ itself is supported.
The Grothendieck theorem, with the orientation just checked, now
applies to the contractible nonempty fibres. For a universal sequence,
the supported stages are exactly a nonempty terminal tail of $\N$;
the same argument applies. Empty target type causes no exception:
every relative path in its one-object identity category is unique.
The independent program \texttt{a\_f\_interval\_prefixmax.py} checks
all $961$ pairs of types of length at most four: $1756$ relative
endpoint functors, including $285$ nonmonotone ones, $9219$ arbitrary
transformations and $7208$ marked transformations. All stated
functoriality and marked-component assertions pass. A further $4565$
checks verify compatibility with source type pullback in lengths at
most three. These finite checks corroborate the general segment proof.
\end{proof}

\begin{theorem}[The missing C.3/C.7 comparisons, with specified maps]
\label{F:unmarked-path-comparison}
For every small $(M,\Sigma)$ and every pair of endpoints the inclusion
\[
 B\mathcal W(M;a,b)\longrightarrow B\mathcal A(M;a,b)
\]
is a weak equivalence. For every universally cofinal actual projective
sequence $t$, both maps in the specified span
\[
 B\colim_n F_{t_n}(a,b)\ \longleftarrow\
 B\mathcal E_t(M;a,b)\ \longrightarrow\
 B\mathcal A(M;a,b)
\]
are weak equivalences. These give the asserted localization mapping
space and its compatibility with composition. In particular the
path-model prerequisites of the supplied C.28 are valid with arbitrary
interior vertical arrows, under the exact marking convention of C.20.
\end{theorem}
\begin{proof}
\textbf{A universal diagram for one simplex.}
Let $\sigma=(p_0\to\cdots\to p_r)$ be a simplex in
$N\mathcal A(M;a,b)$. Write its type maps as $u_{ij}:I_{\tau_i}\to
I_{\tau_j}$ and its transformations as $\eta_{ij}:p_i\Rightarrow
p_ju_{ij}$. Form the covariant Grothendieck construction of the
interval diagram over $[r]$. Its objects are $(i,k)$; an arrow to
$(j,l)$, for $i\leq j$, is an arrow $u_{ij}(k)\to l$ in
$I_{\tau_j}$. Composition applies $u_{jv}$ and then composes in
$I_{\tau_v}$. This formula proves the category laws.

Identify the left endpoint column with one object $A$ and the right
endpoint column with one object $B$, and set their vertical arrows to
identities. Take this quotient as a categorical presentation, denoted
$G_\sigma$. It is small, although the quotient can have infinitely
many arrows. Evaluation at the given diagram is a strict functor
$G_\sigma\to M$: the defining equations are precisely naturality
and composition for the $\eta_{ij}$.

Mark the negative row arrows and, for every pair $i\leq j$, every
cocartesian arrow $(i,k)\to(j,u_{ij}k)$ whose evaluated component
$\eta_{ij}(k)$ belongs to $\Sigma$; close under identities and
composition. This all-pairs condition is essential. It makes
$G_\theta\to G_\sigma$ relative for every face $\theta\subseteq
\sigma$, including a face which omits intermediate rows. The functors
for iterated faces agree strictly. Evaluation remains relative.

Projection to the last row gives relative functors
\[
 r:G_\sigma\longrightarrow I_{\tau_r},\qquad
 s:I_{\tau_r}\longrightarrow G_\sigma,
 \qquad rs=1,\qquad 1\Rightarrow sr.
\]
On $(i,k)$, $r$ has value $u_{ir}k$. Cocartesian arrows give the last
transformation. It descends through the endpoint quotient because
its components there are identities. Every additionally marked
cocartesian arrow maps to an identity under $r$, so both functors
are indeed relative. If the last interval is empty, its two endpoints
coincide and the quotient also identifies $A=B$; the same formulas
apply without choosing two different objects.

\textbf{Contractibility of the three carriers.}
Postcomposition with $r,s$ and their transformation gives categorical
homotopy equivalences on $\mathcal A$ and on every $\mathcal E_t$.
The previous lemma therefore makes
$B\mathcal A(G_\sigma;A,B)$ and
$B\mathcal E_t(G_\sigma;A,B)$ contractible.

For $\mathcal W$ use Lemma~\ref{F:marked-hammocks}, not an
unmarked-vertical assertion. An arbitrary natural transformation
between relative functors induces a natural transformation of their
infinity localizations. To check this directly, functors
$C\times[1]\to D$ are functors $C\to\operatorname{Fun}([1],D)$;
an arrow in the latter functor category is an equivalence exactly
when both endpoint arrows are. The localization universal property
therefore identifies localization of $C\times[1]$ at the copies of
the marking with $LC\times[1]$. If the original transformation is
the identity at $A,B$, its naturality gives a homotopy between the
two induced maps of $\operatorname{Map}(A,B)$. Applying this to
$1\Rightarrow sr$ proves that $r,s$ are inverse on these mapping
spaces; no interior component is being inverted by this argument.

The localization of $I_\tau$ at its negative edges is equivalent to
the ordinal with one edge for each positive edge of $\tau$. Indeed,
functors from the free oriented line are choices of arrows on its
edges; requiring the negative ones to be equivalences contracts each
negative block, with contractible choices of inverse and composite.
This verifies the universal property, including higher choices.
The mapping space between its two endpoints is contractible.
Consequently $B\mathcal W(G_\sigma;A,B)$ is contractible too.

\textbf{The relative carrier argument in every degree.}
Here are the details that a one-square argument does not supply.
Represent a finite sphere, disk, or homotopy diagram by a finite
simplicial complex $K\to N\mathcal A(M;a,b)$ after subdivision and
relative simplicial approximation. Let $P(K)$ be its face poset.
The construction above gives strict covariant diagrams over $P(K)$
of the categories $\mathcal W(G_\sigma)$ and
$\mathcal A(G_\sigma)$. Their classifying spaces are objectwise
contractible. Hence the induced map of their covariant Grothendieck
constructions is a weak equivalence, and each projection to $P(K)$
is a weak equivalence.

The latter construction for $\mathcal A$ has a strict section: at
$\sigma$ choose its last-row path. For a face inclusion
$\theta\subseteq\sigma$, use the universal arrow from the last
row of $\theta$ to the last row of $\sigma$. Composition is the
defining composition of the universal diagram. After evaluation in
$M$ this section is exactly the last-vertex map of the barycentric
subdivision of the original simplex diagram.

If a subcomplex $L\subseteq K$ already maps to
$N\mathcal W(M;a,b)$, every corresponding cocartesian component is
marked by construction. The very same section on $P(L)$ therefore
lifts to the marked Grothendieck construction. Replace its weak
equivalence to the arbitrary-vertical construction by the mapping-path
fibration. It is a trivial fibration. The inclusion
$|P(L)|\subseteq|P(K)|$ is a CW cofibration, so the section lifts
relative to this prescribed boundary. Retracting the mapping-path
replacement yields a lift up to a homotopy fixed on the boundary.
Evaluation in $M$ proves vanishing of every relative homotopy group,
including the path and component assertions. This proves the first
displayed weak equivalence.

Use exactly the same argument with $\mathcal E_t(G_\sigma)$ in
place of $\mathcal W(G_\sigma)$. On a boundary simplex coming from
$\mathcal E_t(M)$ retain its given stage at each row. These stages
decrease along its arrows. For a face inclusion the last-row stages
and the given sequence transition consequently give the required
arrow in $\mathcal E_t(G_\sigma)$; all face compositions still agree.
Thus the boundary section again lifts strictly. Its objectwise
contractible carriers prove that the specified
$B\mathcal E_t(M)\to B\mathcal A(M)$ is a weak equivalence.
No assertion that $t$ is cofinal for arbitrary presheaves is used;
Proposition~\ref{F:C23-uniform} explains why that assertion would
be false.

\textbf{The actual colimit and composition.}
Write $K_n=F_{t_n}(a,b)$, with its actual pullback transitions
$K_m\to K_n$ for $m\leq n$. Sending a stage object to its colimit
class defines the left-hand functor in the stated span: an arrow
$p\to t(n\to m)^*q$ becomes an arrow from $[p]$ to $[q]$.
Any finite diagram in $\mathcal E_t$ pulls back to a common later
stage $N$. Its pullback has a canonical transformation to the original
diagram in the fibred total category, given by identity fibre arrows.
Conversely every finite list of colimit objects, arrows and equations
occurs at a common stage. Applying these two facts to finite sphere
and disk diagrams proves that this specified functor induces all
homotopy-group isomorphisms. It is the fibred version of the filtered
comparison in Proposition~\ref{F:C22-types}.

Finally concatenation is a strict functor on the all-type categories
and commutes with their inclusion $\mathcal W\to\mathcal A$.
The marked comparison identifies it with localized composition by
Lemma~\ref{F:marked-hammocks}. The sequence span therefore transports
the same composition. Where block concatenation gives an actual
stage operation, as for powers of a fixed mixed type, its square with
the all-type functor commutes strictly. Otherwise a chosen later
type receives the concatenated finite data, and the type-map
contraction makes the induced operation agree with this transported
one. Thus the compatibility is with the specified comparisons, not
only an abstract equivalence of individual mapping spaces.
\end{proof}

\IfFileExists{sections/F_global.tex}{\subsection{The global hypothesis in C.28: checking the prerequisite rather than assuming it}
\begin{lemma}[The no-small-quotients resolution used by INDLOCAL]\label{F:large-resolution}
For a connected CW complex $W$ and an infinite cardinal $\lambda$ at
least its number of cells, there is an acyclic map $t:BG\to W$ with
$|G|\leq\lambda$ such that every homomorphism from $G$ to a group of
cardinality less than $\lambda$ kills $\ker\pi_1(t)$. Here acyclic means
an isomorphism in homology with every local coefficient system on $W$.
\end{lemma}
\begin{proof}
We check the supplied \textsc{OWN\_RESULTS} Theorem~3.1, including its
previously unopened input. Leary's Theorem~A \cite[p.~2]{leary2013}
gives an aspherical locally CAT(0) complex mapping to any simplicial
complex by a homology equivalence for all local systems. A CW complex
has a simplicial model of cardinality at most the maximum of its cell
cardinality and $\aleph_0$. Leary's construction is finite on each finite
subcomplex and commutes with their unions. It therefore has that same
cardinal bound, and its fundamental group does too. This supplies the
exact Kan--Thurston input in the successor step; the original 1976 paper
need not be accepted from memory.

\textbf{The large simple acyclic overgroup input.} The embedding
construction checked in Theorem~\ref{A:acyclic} gives an acyclic
overgroup of any group without increasing an infinite cardinal bound.
For completeness the simple-envelope step used in the supplied argument
can be obtained as follows. For every $g\ne1$ and $h$ in a group $L$,
introduce a free letter $t=t_{g,h}$ and put
\[
 w=(gtg^{-1}t^{-1})(gt^2g^{-1}t^{-2}).
\]
Both $w$ and $wh$ have infinite order in $L*F(t_{g,h})$: their cyclically
reduced normal forms have positive alternating length (length eight,
or six after the possible boundary cancellation). Adjoin an HNN letter
$s$ with $s^{-1}ws=wh$. The cyclic subgroups being identified are
infinite cyclic, so the normal-form theorem embeds the old group.
Moreover $h=w^{-1}s^{-1}ws$ belongs to the normal closure of $g$.
Iterating over all pairs and then over countably many stages produces
an overgroup in which every nonidentity element normally generates the
whole group. Start with a free product with a free group on $\lambda$
letters to ensure cardinality exactly $\lambda$. Alternately take this
simple envelope and an acyclic envelope. In the union, every homology
class dies at a later acyclic stage, and every nonidentity element
normally generates every later stage. Thus the union is simple,
acyclic, of cardinality $\lambda$, and contains the specified group.
This checks the exact strength of CATHOT B.7 needed here.

\textbf{Start of the cellular induction.} Use a single vertex for $W$
and take $G_1$ to be the free group of its one-skeleton. Then
$BG_1\simeq W^{(1)}$, with zero kernel. Well-order higher cells so their
boundaries occur earlier; at limit stages use unions.

\textbf{The comparison for one attached cell.} Suppose $t_\alpha:BG_\alpha
\to W_\alpha$ has been constructed and the next cell has attaching map
$a:S^m\to W_\alpha$, $m\geq1$. For $m=1$, take $H'=\Z$ and lift the
attaching loop using the surjection on fundamental groups. For $m\geq2$,
form $E=S^m\times^h_{W_\alpha}BG_\alpha$. Acyclicity is preserved by
homotopy pullback, so $E\to S^m$ is acyclic and $E$ is connected.
Take the preceding bounded Kan--Thurston resolution $BH'\to E$.
The projection to $BG_\alpha$ is represented by a homomorphism
$u:H'\to G_\alpha$, and there is a specified homotopy
$t_\alpha Bu\simeq aq$ with $q:BH'\to S^m$ acyclic. The pullback can be
modeled with at most $\lambda$ simplices: finite words, horns, and
countably many simplicial replacement steps preserve an infinite
cardinal bound. Hence $|H'|\leq\lambda$.

\textbf{The amalgam and its map.} Embed $H'$ in two simple acyclic
groups $A',A''$ of size $\lambda$ and set
\[
 G_{\alpha+1}=(G_\alpha\times A')*_{H'}A'',
 \qquad h\longmapsto(u(h),i'(h)),\quad h\longmapsto i''(h).
\]
Both embeddings are injective, so both factors and $G_\alpha$ embed.
The Bass--Serre tree shows that the double mapping cylinder of their
classifying spaces is $BG_{\alpha+1}$: its universal cover is the tree
of contractible vertex and edge spaces. Map the first vertex space by
$t_\alpha\operatorname{pr}_1$, the second to the center of the new cell,
and the intervening cylinders by the chosen homotopy and the radial
cone on $q$. This defines $t_{\alpha+1}$ with the correct restriction.

\textbf{Acyclicity and the quotient property.} With a local system from
$W_{\alpha+1}$, the $A'$ factor acts trivially and is acyclic. The product
projection is consequently a local-coefficient homology equivalence;
the other two pieces are handled by acyclicity of $q$ and $A''$.
Mayer--Vietoris gives acyclicity of $t_{\alpha+1}$. A homomorphism into a
group smaller than $\lambda$ kills both simple factors. The amalgamation
relation then kills $u(H')$. By induction it factors through
$\pi_1W_\alpha$, and it kills the attaching loop when $m=1$; for $m\geq2$
there is no new fundamental-group relation. Van Kampen therefore makes
it factor through $\pi_1W_{\alpha+1}$. These are the comparison maps
needed in the supplied proof, not merely an equality of homology groups.

\textbf{Limits.} At a limit ordinal every finite chain and every group
word occurs earlier. Local-coefficient homology and the presentation of
the fundamental group commute with these unions. The factor maps are
unique because $G_\alpha\to\pi_1W_\alpha$ is onto, so they assemble.
The number of stages and every stage size are at most $\lambda$.
This completes all induction and cardinal steps of the asserted resolution.
\end{proof}

\begin{theorem}[The global ind-groupoid comparison used in C.28]\label{F:global-ind-groupoids}
Let $K:I\to\Cat_{\mathrm{hot}}$ be a small filtered diagram. If
$H(X)=\colim_i\pi_0\operatorname{Fun}(X,K_i)$ inverts every nerve
weak equivalence between small categories, then the canonical
ind-morphism $(K_i)\to(\Pi_1K_i)$ is an isomorphism. The global
quantifier on tests is essential.
\end{theorem}
\begin{proof}
\textbf{Descent and its specified map.} Invariance descends $H$ to a
presheaf $F$ on $\mathrm{Hot}$. For any presheaf $T$ there, Yoneda gives
\[
 \operatorname{Nat}(F,T)
 =\operatorname{Nat}(H,T\circ B)
 =\lim_i T(BK_i)
 =\operatorname{Nat}\bigl(\colim_i[-,BK_i],T\bigr).
\]
Thus $F(Y)=\colim_i[Y,BK_i]$, and the identification sends a functor to
its classifying map. The localization fact used here follows by replacing
a simplicial set by its category of simplices: its nerve realizes to the
original space, the last-vertex functor compares it to an original nerve,
and these natural equivalences represent maps and homotopies by roofs.
This also checks the naturality, which the subsequent cancellation needs.

\textbf{The two comparisons.} Choose an infinite cardinal $\mu$ bounding
all arrow sets of the stages, and choose $\lambda>\mu$ large enough for
a test space $Y$. Componentwise apply Lemma~\ref{F:large-resolution} to
obtain $t:T=\coprod BG_c\to Y$ with $p:G_c\twoheadrightarrow\pi_1Y_c$.
Every functor $\mathbf BG_c\to K_i$ lands in the core groupoid of $K_i$.
Its image automorphism group has at most $\mu$ elements, so the functor
factors uniquely through $\mathbf B\pi_1Y_c$. The same assertion for a
natural transformation follows by choosing a lift of each group element
under $p$ in the naturality equation. Hence
\[
 \operatorname{Fun}\bigl(\coprod\mathbf B\pi_1Y_c,K_i\bigr)
 \cong \operatorname{Fun}\bigl(\coprod\mathbf BG_c,K_i\bigr).
\]
This is the actual functor-category comparison in INDLOCAL B.1.

\textbf{The injectivity needed for cancellation.} An acyclic map
$t:T\to Y$ is onto on fundamental groups. The fold
$Y\cup_T^hY\to Y$ is an isomorphism on $\pi_1$ by van Kampen and on
homology with every local coefficient system by Mayer--Vietoris.
Using coefficients $\Z[\pi_1Y]$ gives a homology equivalence of simply
connected universal covers, hence a homotopy equivalence. Its two
sections are equal in $\mathrm{Hot}$. It follows that
$[Y,Z]\to[T,Z]$ is injective for every space $Z$. Filtered colimits
preserve this injectivity: an equality is witnessed at a common stage,
where the preceding injection applies.

\textbf{Cancellation and the unit.} Let $v:Y\to P_1Y=\coprod B\pi_1Y_c$.
The composites $vt$ and $Bp$ agree because their maps on fundamental
groupoids agree and the target is a $1$-type. The functor-category
comparison proves $F(vt)$ bijective. The previous paragraph proves
$F(t)$ injective. Since $F(vt)=F(t)F(v)$, surjectivity of the composite
also makes $F(t)$ surjective; therefore both $F(t)$ and $F(v)$ are
bijections. In particular $H$ inverts every categorical unit
$X\to\Pi_1X$. Yoneda now identifies the unit of the induced reflection
on ind-objects with $(K_i)\to(\Pi_1K_i)$: left Kan extension of each
representable is represented by its groupoid completion, and the unit
on their colimit is invertible exactly when $H$ inverts these units.
This verifies the filtered proof of INDLOCAL B.1 used in C.28; it does
not apply this theorem to a presheaf invariant only on finite tests.
\end{proof}

\begin{proposition}[C.28: the finite-test comparison under global invariance]
\label{F:C28}
Suppose a sequence $(K_n)$ has globally invariant ind-presheaf $H$ as
above. Put $Z=\hocolim_n BK_n$ and
$\Gamma=\colim_n\Pi_1K_n$. The canonical map $Z\to B\Gamma$ is a
homotopy equivalence. For every finite category $Y$ the specified map
\[
 c_Y:H(Y)\longrightarrow[BY,Z]
\]
is a bijection. The same conclusion holds when $\Pi_1Y$ is finitely
presented. This is the valid finite-category formula at VII
ms.~129--132, under the global invariance premise of that paragraph.
\end{proposition}
\begin{proof}
\textbf{The map of spaces.} The stagewise units give the displayed
canonical map. Finite edge words, relations, and basepoint paths occur
at a common stage, so it is bijective on components and an isomorphism
on every fundamental group. The preceding theorem says $H$ depends
only on $\Pi_1$. A finite face poset for $S^r$, $r\geq2$, has trivial
fundamental groupoid, so its value equals the value at a point.
The descended formula in the proof above and compactness of the finite
sphere identify this with $[S^r,Z]=[*,Z]$. In each component this says
that every based homotopy class lies in the orbit of zero; that orbit is
just zero. Thus all higher homotopy groups vanish, and the canonical map
between these $1$-types is a homotopy equivalence.

\textbf{The finite-presentation step.} The groupoid $\Pi_1Y$ has a finite
presentation from objects, arrows, identity equations, composition
equations, and formal inverses. A functor to $\Gamma$ and a natural
isomorphism between two such functors are finite data with finitely many
equations. They therefore occur at a common stage. Consequently
\[
 H(Y)\cong\colim_n\pi_0\operatorname{Fun}(Y,\Pi_1K_n)
 \cong\pi_0\operatorname{Fun}(Y,\Gamma)
 \cong[BY,B\Gamma]\cong[BY,Z].
\]
Every arrow in this chain is induced by the stagewise units or their
classifying maps. Hence this is $c_Y$, not an unspecified bijection.
This checks both (C.28.1) and (C.28.2) and explains why no compactness
assumption about the whole nerve of $Y$ is needed.

\textbf{Why arbitrary tests are different.} Even the discrete groupoid
sequence $\{1,\ldots,n\}$ has a globally invariant ind-presheaf.
On an infinite discrete test set its colimit consists of finite-image
functions, whereas maps to its union $\N$ include functions with infinite
image. This is a fully explicit failure of the unrestricted version of
$c_Y$, with the same comparison map. Thus the qualification in C.28 is
necessary.
\end{proof}

\begin{lemma}[The required diagram-localization sphere test]
\label{F:diagram-sphere}
Let $\mathcal M$ be bounded finite-dimensional rational chain complexes,
marked at quasi-isomorphisms, and let $P$ be a finite poset. For constant
diagrams $A=\Q[0]$, $B=\Q[2]$, there is a specified isomorphism
\[
 \Hom_{\mathcal M^P[\Sigma_P^{-1}]}(A_P,B_P)
       \cong H^2(BP;\Q).
\]
In particular the values for a finite face poset of $S^2$ and for a
point are $\Q$ and the zero group, respectively.
\end{lemma}
\begin{proof}
Use the diagram of chain complexes
$Q_P(x)=C_*(N(P\downarrow x);\Q)$. Its augmentation to $A_P$ is an
objectwise quasi-isomorphism because each slice has terminal object
$x$. In degree $r$ it is a finite sum of projective representable
diagrams, one for each strict chain $x_0<\cdots<x_r$, with summand
$\Q[P(x_r,-)]$. Thus $Q_P$ is a bounded complex of finite projective
diagrams. The category of diagrams of bounded finite-dimensional
complexes on finite $P$ is exactly the category of bounded complexes
of finite-dimensional rational diagrams on $P$: finiteness supplies
a common bound on their degrees.

For completeness, a bounded projective complex has acyclic Hom complex
into every acyclic diagram complex. This follows by induction on its
number of nonzero terms, starting with a shifted projective and using
the short exact sequence separating its top term. Consequently it
computes morphisms in localization by quasi-isomorphisms: the cone of
each inverted map is acyclic, so applying Hom gives a bijection on
chain-homotopy classes; roofs then have unique such representatives.
The desired morphisms are therefore chain-homotopy classes of maps
$Q_P\to B_P$.

The colimit of each representable summand is one copy of $\Q$, and
its boundary maps are exactly deletion of chain vertices. Hence
$\colim_P Q_P=C_*(NP;\Q)$, with its actual augmentation. The
colimit--constant-diagram adjunction, on chain maps and on chain
homotopies, now identifies the preceding classes with
$H^2(\operatorname{Hom}(C_*(NP;\Q),\Q))=H^2(BP;\Q)$.
This is a calculation in the strict diagram localization itself;
it assumes no general comparison with a path-space model.
\end{proof}

\begin{proposition}[C.26: the loop formulas do not imply global invariance]
\label{F:C26-complexes}
The countable category of bounded finite-dimensional rational chain
complexes, marked at quasi-isomorphisms, satisfies both pointed loop and
suspension formulas but fails the global invariance condition of IX.
\end{proposition}
\begin{proof}
\textbf{The precise model hypotheses.} Use finite rational matrices to
choose a countable small model, with degreewise injections as
cofibrations and degreewise surjections as fibrations. It has all finite
limits and colimits, the exact hypothesis M0 of Dwyer--Kan
\cite[Section~1.4(iv)]{dwyerkan1980function}; arbitrary coproducts are
not required by that source. An acyclic vector-space complex is
contractible by splitting its cycles. In a lifting square, a graded lift
has a differential defect in the Hom complex from the cokernel of the
injection to the kernel of the surjection. If either given map is a
quasi-isomorphism, that Hom complex is contractible, and subtracting a
boundary corrects the defect. This proves the lifting axiom.

\textbf{Factorizations and base change.} Write $C(A)=\operatorname{Cone}(1_A)$.
The maps $A\to B\oplus C(A)\to B$, the first $(f,j)$ and the second
projection, are a cofibration followed by a trivial fibration. The maps
$A\to A\oplus C(B)[-1]\to B$, the first inclusion and the second $(f,q)$,
are a trivial cofibration followed by a fibration. Kernels, cokernels,
and exactness of vector spaces verify all pullback and pushout axioms;
homology gives two-out-of-three. Every object is cofibrant and fibrant.

\textbf{The mapping spaces.} The cosimplicial complex
$A\otimes N\Q\Delta^\bullet$ augments by quasi-isomorphisms, and its
latching inclusions are the boundary inclusions, hence degreewise
injections. Dwyer--Kan Proposition~4.4 applies, with fibrant target $B$,
and identifies the derived mapping space with the simplicial abelian
group $r\mapsto\Hom_{\mathrm{Ch}}(A\otimes N\Q\Delta^r,B)$.
Its normalized complex is the good truncation
$\tau_{\geq0}\operatorname{Hom}_*(A,B)$: the value on the top simplex is
the normalized generator and its boundary is the alternating face sum;
in degree zero one takes cycles. Thus its $i$th homotopy group is
$H_i\operatorname{Hom}_*(A,B)$. This verifies the actual comparison used
in C.26, including its finite-limit hypothesis.

\textbf{The formulas and their failure to imply invariance.} Shifts
$\Sigma A=A[1]$ and $\Omega B=B[-1]$ give the based loops of these
mapping spaces. Equivalently the displayed cone factorizations compute
$0\amalg_A^h0$ and $0\times_B^h0$; mapping-space universal properties
supply their coherent group laws. For $A=\Q[0]$ and $B=\Q[2]$ the mapping
space is $K(\Q,2)$. Proposition~\ref{F:C20-paths} represents the actual
constant-endpoint diagram-localization functor by the indicated
ind-path object, on all small tests. If it were globally invariant,
Theorem~\ref{F:global-ind-groupoids} would force its value on a face
poset of $S^2$ to equal its value on a point, since both have terminal
fundamental groupoid. Lemma~\ref{F:diagram-sphere} computes those
actual values as $\Q$ and a singleton. This is impossible. This repairs
the dependency on the unavailable C.3/C.7 path-space proof in the
supplied final paragraph, while confirming the claimed counterexample
and every model-category calculation that precedes it.
\end{proof}
}{}
\subsection{C.28 at V ms.~46--49: stages, unions, and strict retractions}
Call a category $A$ local when the specified comparison
$\pi_0\operatorname{Fun}(X,A)\to[BX,BA]$ is bijective for every small
$X$. This is stronger than Thomason fibrancy. The model structure used
here is the one stated and proved in
\cite[Theorem~2.3]{may2017posets}: its generating trivial cofibrations
are $c\operatorname{Sd}^2\Lambda_i^r\to c\operatorname{Sd}^2\Delta^r$,
and fibrancy means that $\operatorname{Ex}^2NA$ is Kan. These generating
maps have finite source and target and are full sieve inclusions: the
twice-subdivided simplex is the poset of chains of faces, and a
subcomplex gives the downward-closed subposet of its face chains.
These are exactly the hypotheses used in the following checks.

\begin{proposition}[The positive presentation assertions in C.28]
\label{F:C28-presentations}
A globally invariant sequential ind-object admits a presentation by
full inclusions of local Thomason-fibrant categories. The union is
Thomason-fibrant. It need not be local or represent the ind-object.
Filtered colimits of Thomason-fibrant categories are fibrant, and
objectwise weak equivalences of filtered diagrams induce weak
equivalences of their colimits.
\end{proposition}
\begin{proof}
\textbf{The groupoid presentation.} Theorem~\ref{F:global-ind-groupoids}
gives the specified ind-isomorphism to the actual groupoid sequence
$\Gamma_n=\Pi_1K_n$. Let $D_n$ be its covariant Grothendieck construction
over $\{0<\cdots<n\}$. Its objects are $(i,x)$ with $x\in\Gamma_i$;
for $i\leq j$ its arrows $(i,x)\to(j,y)$ are the arrows
$\Gamma_{ij}(x)\to y$ of $\Gamma_j$. The product of labels is
$h\,\Gamma_{jk}(g)$, which verifies associativity and identities.
The terminal-fibre retraction $r_n:D_n\to\Gamma_n$ sends objects and
labels forward. Its section $j_n$ satisfies $r_nj_n=1$, and identity
labels give $1\Rightarrow j_nr_n$. Thus $D_n$ is categorically homotopy
equivalent to a groupoid and is local: functors to a groupoid factor
through fundamental groupoids, and their transformations are precisely
the natural isomorphisms classifying unpointed maps of classifying
spaces. Applying the two inverse zigzags transports this comparison to
$D_n$.

\textbf{Fullness and fibrancy.} The $D_n$ are full subcategories of
$D_{n+1}$, with strictly compatible retractions. Given a generating
horn $A\subset B$ and $v:A\to D_n$, first extend $r_nv$ to
$w:B\to\Gamma_n$. A groupoid has Kan nerve and hence is
Thomason-fibrant, so this extension exists. Retain $v$ on old objects
and send each new object to its value under $w$ in the terminal fibre.
Retain old arrows and label all new arrows by $w$. Since $A$ is a full
sieve, there is no arrow from a new object to an old one; all other
new arrows have allowed terminal-fibre targets. Mixed compositions
are exactly multiplication of their labels in $\Gamma_n$, so the
extension is a strict functor. This checks the lifting assertion in
the written proof of C.28.

\textbf{Filtered passages.} A lifting problem for a generating horn
in a filtered colimit is finite data and finite equations, so it occurs
at a stage and is solved there. This proves fibrancy of the colimit.
Nerves commute with filtered colimits by the same finite-data argument;
their ordinary colimits agree weakly with homotopy colimits as checked
in Proposition~\ref{F:C22-types}. Objectwise weak equivalences therefore
induce weak equivalences of these colimits. For an arbitrary actual
sequence, compatible fibrant replacements also exist: choose trivial
cofibrations $K_n\to R_n$ to fibrant targets, and extend
$K_n\to K_{n+1}\to R_{n+1}$ across $R_n$ by lifting. A sequence is
freely generated by its adjacent arrows, so no coherence equations
remain to be solved. None of these lifting arguments asserts locality
of the union. The next counterexample verifies its failure.
\end{proof}

The previously missing all-type path-space comparison is now proved
in Theorem~\ref{F:unmarked-path-comparison}. Its relative universal
diagrams supply the coherence in every homotopy degree and check the
specified sequential span. The independent C.20 component proof,
C.26 diagram test and C.28 finite-test theorem do not need that repair.

\begin{proposition}[The union and uniqueness counterexample in C.24/C.28]
\label{F:C28-union-counterexample}
There is a sequence of finite local fibrant categories with full
inclusions whose fibrant union $T$ is not local. It has an endofunctor
$s$ equal to the identity as an ind-morphism in
$\operatorname{Ind}(\Cat_{\mathrm{hot}})$ but not a nerve weak
equivalence. Thus ind-isomorphisms do not give the asserted uniqueness
of unions at V ms.~46--49.
\end{proposition}
\begin{proof}
\textbf{The tables and stage retractions.} Let
$G_n=\operatorname{Sym}\{0,\ldots,n\}$ and take $T_n$ with objects
$0,\ldots,n$, with $\Hom(i,j)=G_j$ for $i\leq j$, and with no arrows
otherwise. Compose labels by
$(h:j\to k)(g:i\to j)=h\,\iota_{jk}(g)$.
Group multiplication verifies every category equation. This is the
preceding $D_n$ for inclusions of one-object groupoids, so its specified
terminal-fibre section, retraction and identity-label transformation
make it local and fibrant; the transitions are full inclusions.

\textbf{The actual ind-homotopies.} Set $s(i)=i+1$ and shift the support
of every permutation one place, fixing zero. This commutes with label
composition. For fixed $n$, both the inclusion and the shift
$T_n\to T_{n+1}$ have identity-label transformations to their
corresponding functors constant on objects at $n+1$. The latter two
are naturally isomorphic by the cycle $c=(0\ 1\ \cdots\ n+1)$,
because $cgc^{-1}=s(g)$ for every $g\in G_n$.
The resulting three-arrow zigzag proves equality with the identity
ind-morphism. This is an actual zigzag at the next stage, not just an
equality of maps of classifying spaces.

\textbf{Failure on the union.} The stage retractions assemble to
$r:T\to\mathbf BG$, where
$G=\operatorname{Sym}_{\mathrm{fin}}(\N)$. The preceding filtered
comparison proves $Br$ a weak equivalence. Under this identification,
$Bs$ induces the group shift. Its image fixes zero and omits $(0\ 1)$,
so it is not a fundamental-group isomorphism and cannot be homotopic
to the identity. A functor $\mathbf BG\to T$ selects one object $i$
and has all its arrow images in the finite group $G_i$. Its classifying
map cannot invert $Br$, even up to unpointed homotopy, since conjugation
preserves finiteness of the image. Thus the local comparison for this
source fails. The weak equivalence $r$ between fibrant categories cannot
have a $\Cat_{\mathrm{hot}}$ inverse either.

\textbf{Finite verification and the other printed obstructions.}
The independent program \texttt{f\_union\_models.py} supplies all
tables through $T_3$, with arrow counts $1,5,23,119$. It checks every
identity and composable triple (respectively $1,23,917,73037$), each
retraction and all three transformations into the next stage.
For the separate unbounded-length obstruction, finite intervals of the
infinite alternating fence contract by deleting endpoints. Its full
union admits no finite categorical homotopy from identity to a constant:
a zigzag of length $l$ would move every vertex by undirected distance
at most $l$, while distances to a fixed vertex are unbounded. Finally
the terminal fork $0\to1\leftarrow2$ is fibrant: a functor on a full
sieve extends by the terminal object on new objects. Its vertices $0,2$
are connected but have no arrow in either direction. This separately
refutes the one-transformation inference.
\end{proof}

\begin{proposition}[The representation and retraction qualifications in C.28]
\label{F:C28-retractions}
For full-inclusion sequences, an ind-morphism in $\Cat$ is exactly a
functor on unions whose restriction to each source stage lands in some
target stage. Equality in $\operatorname{Ind}(\Cat_{\mathrm{hot}})$
does not imply equality or homotopy of the induced union functors.
A globally invariant ind-object is represented in $\Cat_{\mathrm{hot}}$
exactly when it is isomorphic to a constant ind-object.

If $A$ is Thomason-fibrant, $i:A\hookrightarrow B$ is a Thomason
cofibration, and $p:B\to A$ has $pi$ categorically homotopic to $1_A$,
there is a strict retraction $r$ in the same ordinary homotopy class as
$p$. If $A$ is also local, the equality holds in
$\Cat_{\mathrm{hot}}$. For a groupoid target the latter assertion
holds for any full inclusion injective on objects.
\end{proposition}
\begin{proof}
\textbf{Representation and bounded union maps.} The formula for these
morphisms is
\[
 \Hom_{\operatorname{Ind}(\Cat)}((A_n),(B_m))
    =\lim_n\colim_m\Hom(A_n,B_m).
\]
With literal full inclusions, equality
in the colimit means equality of the corresponding union maps. Compatible
representatives therefore glue to one functor on the union; conversely
a functor satisfying the displayed boundedness yields those compatible
classes. Replacing Hom by categorical homotopy classes only provides
eventual zigzags. Proposition~\ref{F:C28-union-counterexample} verifies
their failure to glue even to an ordinary homotopy of union maps.
The constant ind-embedding is fully faithful, and its associated
presheaf is the representable one. Yoneda gives the claimed exact
representation criterion. A weak equivalence between local categories
has a categorical-homotopy inverse: take its inverse in $\mathrm{Hot}$,
represent it using locality of its target, and use injectivity of both
local comparisons for the two composite equations. The preceding $r$
is an explicit failure when only fibrancy is assumed.

\textbf{The path-object construction, checking INDLOCAL D.1.}
Represent $pi\simeq1_A$ by a functor $A\times Z\to A$ for a finite
oriented interval $Z$ with endpoints $0,z$. The constant-path functor
$c:A\to\operatorname{Fun}(Z,A)$ is a categorical homotopy equivalence:
successive endpoint deletions contract $Z$, and their transformations
induce the inverse zigzags on functor categories. Factor endpoint
evaluation as a trivial cofibration followed by a fibration,
\[
 \operatorname{Fun}(Z,A)\xrightarrow{j}P
       \xrightarrow{(d_0,d_1)}A\times A.
\]
Since $A$ is fibrant, both projections of $A\times A$ are fibrations,
so $d_0,d_1$ are fibrations. The map $jc$ is a weak equivalence and
$d_0jc=d_1jc=1$, so both are trivial fibrations. The adjoint of the
given homotopy followed by $j$ gives $h:A\to P$, with $d_0h=pi$ and
$d_1h=1$. Lifting against $d_0$ extends it to
$\widetilde h:B\to P$ with $d_0\widetilde h=p$. Put
$r=d_1\widetilde h$. Then $ri=1$ strictly. In $\mathrm{Hot}$ both
$d_0,d_1$ equal the inverse of $jc$, so $[r]=[p]$ there. No finite
natural-transformation zigzag on $B$ follows from this factorization
alone.

\textbf{The two upgrades, checking INDLOCAL D.2.} Locality makes
$\pi_0\operatorname{Fun}(B,A)\to[BB,BA]$ injective and supplies the
first upgrade. For an actual groupoid, compose the given zigzag into
an isomorphism $\alpha:pi\Rightarrow1_A$. Identify $A$ with its full
image. Put $r(b)=b$ on $A$ and $r(b)=p(b)$ outside; use
$\beta_b=\alpha_b$ on $A$ and identity outside. On an arrow $u:b\to b'$
define $r(u)=\beta_{b'}p(u)\beta_b^{-1}$. Cancellation checks composition,
naturality of $\alpha$ gives $ri=1$, and $\beta:p\Rightarrow r$ is the
specified natural isomorphism. These check every qualification of the
retraction sentences used in C.28.
\end{proof}

\subsection{A direct replacement for the missing C.17 example}
\begin{theorem}[Invariant paths with a nonlocal actual union]
\label{F:C17-fence}
Let $C$ have objects $0,1$, identities, and two distinct arrows
$f,g:0\to1$, and mark all its arrows. For every small category $Y$,
the constant-endpoint localization functor is naturally
\[
 H_{0,0}(Y)=
 \Hom_{C^Y[\operatorname{Mor}(C^Y)^{-1}]}(0_Y,0_Y)
 \cong\{v:\pi_0Y\to\mathbb Z\mid v\text{ has finite image}\}.
\]
It is invariant under every nerve weak equivalence. Nevertheless,
the actual symmetric alternating path union is not local and is not
categorically homotopy equivalent to a local category. Its specified
comparison on an infinite discrete test omits functions of infinite
image. This supplies the two C.17 assertions used in C.28 directly.
\end{theorem}
\begin{proof}
\textbf{The cover and the exact all-small-test computation.}
Let $P$ be the integer fence: its only nonidentity arrows are
$2r\to2r-1$ and $2r\to2r+1$. Send an even vertex to $0$, an odd
vertex to $1$, the right-going arrow to $f$ and the left-going arrow
to $g$. This functor $p:P\to C$ has unique lifts of arrows with
either a prescribed source or a prescribed target. The same is
true for nerve simplices, so $Bp$ is a covering map. Its deck
transformations are translations by even integers; $BP$ is the
real line with this cellular subdivision. In particular $BC$ is
the two-edge circle, not merely a space with the same homology.

For an arbitrary, possibly disconnected or empty, $Y$, let
$\operatorname{Fun}_b(Y,P)$ be the full subcategory of functors
whose integer-valued object images are uniformly bounded. Every
finite nerve diagram in it lies in
$\operatorname{Fun}(Y,P_{[-N,N]})$ for some $N$. The finite fence
$P_{[-N,N]}$ contracts categorically to $0$, fixing $0$: delete an
endpoint other than $0$, retract it onto its unique neighbor, and
repeat. At each deletion the retraction is joined to the identity
by a natural transformation in one of the two directions.
Postcomposition gives the same contraction of the functor category,
for this arbitrary $Y$. Consequently
$B\operatorname{Fun}_b(Y,P)$ is weakly contractible. Indeed its
filtered union is the union of these subcategories, and finite
simplicial sphere and homotopy data occur at a common finite stage.

A natural transformation of diagrams in $C$ lifts uniquely after a
lift of either endpoint diagram is chosen. Naturality follows by
uniqueness of the lifted composites in each square. Every component
of such a lift changes its integer height by at most one. It
therefore preserves uniform boundedness. A finite zigzag starting
at the constant $0$ diagram has a bounded lift. Conversely a bounded
lift contracts in a finite fence, so its image is in that same
component of $\operatorname{Fun}(Y,C)$. We have thus obtained a
categorical covering onto exactly the component of $0_Y$.

Two bounded lifts differ by a translation $2v(y)$, constant on
every connected component of $Y$; the function $v$ is uniformly
bounded. Conversely each such function acts freely on the cover,
and these transformations act transitively on every fiber. Its
deck group is therefore the bounded integer functions on $\pi_0Y$.
Bounded and finite-image have the same meaning for a subset of
$\mathbb Z$. The weakly contractible cover computes the fundamental
group at $0_Y$ as this group. Inverting all arrows in a category
gives its fundamental groupoid, which proves the displayed formula.
The identification sends a loop of diagrams to its componentwise
winding number and hence commutes with every restriction in $Y$.
For empty $Y$ both sides are the one-element group. Since weak
equivalences preserve $\pi_0$, this proves global invariance.

\textbf{The actual full-inclusion union.}
Let $I_n$ have vertices $-2n,\ldots,2n$ and its fence orientations,
and let $K_n$ be its endpoint-fixed path category in $C$, with both
endpoints $0$ and arbitrary interior transformations. The type is
$(+-)^{2n}$. Extend paths and transformations by the constant $0$
on both new ends to obtain full inclusions $K_n\subset K_{n+1}$.
The types are universally cofinal: their alternating run counts
are unbounded, or one may allocate each edge of any fixed type to
a matching edge in a sufficiently long alternating word.
Write $U=\bigcup_n K_n$. Lifting a path from its left endpoint
$0\in P$ gives a right endpoint $2w$. This integer $w$ is unchanged
by every endpoint-fixed transformation. The component computation
in Proposition~\ref{F:C20-paths}, applied to the point, identifies
$\pi_0U$ with $\mathbb Z$ by this same winding number.

The winding-zero full subcategory $U_0$ has a particularly literal
description. Its objects are order-preserving maps $h:P\to P$
equal to $0$ outside a finite interval. Its arrows are pointwise
order relations. For each $N$, let $Q_N$ be the full subcategory
of such maps with values in $[-N,N]$. The preceding finite-fence
contraction fixes $0$, so postcomposition contracts $Q_N$ and
preserves the finite-support condition. Thus
$U_0=\bigcup_NQ_N$ is weakly contractible and is connected.

\textbf{The obstruction to a finite categorical homotopy.}
Suppose a zigzag of $l$ natural transformations joined $1_{U_0}$
to a constant functor with value $c$. Evaluate this zigzag at an
object $h$ and then at the fence coordinate $0$. Each arrow in
$P$ changes integer height by at most one, so
\[
 |h(0)-c(0)|\le l
 \qquad\text{for every object }h\text{ of }U_0.
\]
For every positive integer $r$, however, the map
\[
 h_r(j)=\max\{0,\,2r-|j|\}
\]
is such an object and has $h_r(0)=2r$. On its support the parity
of $h_r(j)$ is the parity of $j$, and adjacent values differ by
one; across the support boundary the same fence relation or an
identity applies. These facts check every functor equation and
give the contradiction as $r$ increases.

If $U$ were local, the inclusion $U_0\to U$ and a constant functor
in its component would be joined by a finite zigzag, because
$BU_0$ is contractible. Every functor and transformation in that
zigzag remains in the same component $U_0$, contrary to the last
paragraph. Locality is preserved by categorical homotopy
equivalences, so $U$ cannot be equivalent in that sense to a local
category either. This is the asserted obstruction for the actual
symmetric union, in addition to the separate union counterexample
of Proposition~\ref{F:C28-union-counterexample}.

\textbf{The precise failed comparison.}
For the discrete category $Y=\mathbb N$, the space $BY$ is discrete
and $[BY,BU]=(\pi_0U)^{\mathbb N}=\mathbb Z^{\mathbb N}$.
Evaluation at each point and the winding computation show that
the specified map $H_{0,0}(Y)\to[BY,BU]$ is exactly the inclusion
of finite-image functions. In particular $n\mapsto n$ is omitted.
No mapping-space theorem for unmarked vertical hammocks has been
used in any part of this proof.

The independent program \texttt{f\_fence\_paths.py} records the
complete four-arrow circle table and the covering-star checks,
all relative contractions of fences of radii one through eight,
and the full symmetric path stages $n=0,1,2$. Their object counts
are $1,29,985$, their arrow counts are $1,71,7573$, and all
$1,225,102471$ composable triples respectively satisfy the
category equations. Their connected components have exactly the
computed winding numbers. Sixteen explicit excursion maps are
also checked; the preceding formula proves the obstruction for
every height, independently of this finite range.
\end{proof}

\subsection{The bounded-type and one-sign checks behind C.28}
\begin{proposition}[Checking the projective-sequence classification C.21]
\label{F:C21-sequences}
For an actual projective sequence of types $t_n$, let $r_n$ count
its nonempty constant-sign runs. The following are equivalent:
$r_n$ is unbounded; every type receives a map from some $t_n$;
and the all-small-test formula in Proposition~\ref{F:C20-paths}
holds for every relative category and endpoint pair.
If these conditions fail, the tail of the path diagram is
isomorphic in $\Cat_{\mathrm{hot}}$ to the constant diagram at
one finite alternating type. The identifications commute with
the actual transition maps and localization effects.
\end{proposition}
\begin{proof}
In the allocation description of Lemma~\ref{F:type-contraction},
opposite successive target runs occupy distinct successive source
runs. Thus a type map cannot increase run count, and $r_n$ is
nondecreasing. Unbounded alternating run strings accommodate any
fixed target word, by allocating its successive letters to matching
runs and repeating positions within each run. Conversely, maps to
every alternating word require unbounded run counts. Theorem C.20,
checked in Proposition~\ref{F:C20-paths}, proves the formula in this
case with its actual comparison.

For necessity, if $r_n\le r$, take the finite alternating fence
with $r+2$ vertices, mark every arrow, and choose its two extremal
vertices as endpoints. A constant-sign path run moves by at most
one edge, since no two nonidentity arrows in either direction
compose. Thus all the path fibers are empty. In the groupoid
completion the two endpoints have exactly one arrow between them,
because the nerve is an interval. This is a complete finite
counterexample for every bounded run count, including $r=0$.

Finally bounded $r_n$ stabilizes. At a common positive run count,
a map of types must preserve the first sign: otherwise the target
run pattern needs one additional source run. Let $\sigma$ be this
fixed reduced alternating word. Selecting and repeating within
runs gives maps both ways between each tail type and $\sigma$.
Their composites and identities lie in the same nonempty mapping
posets, which contract by Lemma~\ref{F:type-contraction}.
Consequently their path pullbacks are inverse in
$\Cat_{\mathrm{hot}}$. The two maps to $\sigma$ associated to
each transition lie in that same mapping poset as well. This
checks compatibility with the given sequence, not only separate
equivalences of its stages. Applying
$\pi_0\operatorname{Fun}(X,-)$ preserves these finite uniform
zigzags for every small $X$. The zero-run case is the empty word.
The remaining endpoint-specific criterion is precisely that
every localization word has a $\sigma$ representative and equal
effects have connected $\sigma$ representatives; these are
surjectivity and injectivity of the specified effect map.
\end{proof}

\begin{proposition}[One-sign paths and the exact universal comparison]
\label{F:one-sign-paths}
For any positive length, the all-positive endpoint-fixed path
category is categorically homotopy equivalent to the discrete set
$\Hom_M(a,b)$. The all-negative one is equivalent to the discrete
set $\Sigma(b,a)$. These identifications commute with every
same-sign type pullback after passage to
$\Cat_{\mathrm{hot}}$, and hence hold after applying
$\pi_0\operatorname{Fun}(X,-)$ for every small $X$.

The comparison of positive paths with localization mapping spaces
is an equivalence at all endpoints exactly when every marked arrow
is already an isomorphism in $M$. For negative paths the analogous
all-endpoint condition is that $M$ is a groupoid and
$\Sigma=\operatorname{Mor}M$. At one endpoint pair the exact
condition is that the specified effect map is bijective on
components and every component of the localization mapping space
is contractible. In particular, when all arrows are marked, the
positive comparison holds everywhere exactly for groupoids.
\end{proposition}
\begin{proof}
Compose a positive path to its effect $f:a\to b$. A section sends
$f$ to the path with $f$ first and identities at $b$ afterwards.
The remaining composites from its intermediate objects to $b$
give a natural transformation from the original path to this
section. This is natural in endpoint-fixed ladders because each
square commutes. Thus composition and the section are inverse
up to a single transformation. For a negative path, use its
marked composite $b\to a$, put that arrow first and identities
after it, and use the marked composites from $b$ to each
intermediate object. They give the transformation in the reverse
direction. Closure of the marking ensures all these are paths.
Effect commutes strictly with type pullbacks. The sections then
give their claimed compatibility in $\Cat_{\mathrm{hot}}$.

The positive all-endpoint condition says exactly that the
localization functor $NM\to L(M,\Sigma)$ is fully faithful;
it is already surjective on objects. It is therefore an
equivalence, so it reflects equivalences and every member of
$\Sigma$ was invertible. Conversely localizing already invertible
arrows changes nothing, by the localization universal property.
For negative paths the structural functor
$\Sigma^{\mathrm{op}}\to L(M,\Sigma)$ sends each arrow to an
equivalence. If it is fully faithful, every arrow in $\Sigma$
has its inverse in $\Sigma$. These inverses are also inverses in
$M$, so the preceding argument identifies $L(M,\Sigma)$ with
$NM$. Full faithfulness now says every arrow of $M$ is the inverse
of one of $\Sigma$, proving the asserted necessity. The converse
and the single-endpoint criterion follow directly from the
discrete source mapping sets.

For the finite failure mentioned in C.6 one may use the four-arrow
circle category of Theorem~\ref{F:C17-fence}, with every arrow
marked. Its forward one-edge fiber at $(0,0)$ is a singleton,
whereas its localization is the groupoid completion of the
circle as a space, whose endpoint mapping space has components
$\mathbb Z$. The complete category table was checked there.
This replaces any reliance on the unavailable crown example.
\end{proof}

\paragraph{The composition assertions.}
For any fixed nonempty block $\tau$, paths in $\tau^n$ and
$\tau^m$ concatenate strictly to a path in $\tau^{n+m}$.
Padding by an identity block is a literal functor, so concatenation
defines the ind-operation with its units and associativity.
If the block has both signs, its run counts are unbounded and
Proposition~\ref{F:C21-sequences} identifies its all-small-test
component formula with localization. If it has one sign, the
preceding proposition identifies the operation with ordinary
composition (or inverse composition in the marked subcategory).
Choices of a common type give the same operation because the
mapping-poset contractions act uniformly on paths. These verify
the categorical and ind-operation statements used in C.28.
Their further identification with all higher localization
mapping-space compositions is exactly the remaining comparison
isolated in \ref{F:unmarked-path-comparison}; it is not inferred
from associativity on components.

\subsection{The complete colimit criterion used in C.28}
\begin{theorem}[Step audit of INDLOCAL C.27 and PRETOT F.1/F.4]
\label{F:C27-colimits}
For a small index category $A$, the canonical functor
\[
 \int_A D\longrightarrow\colim_A^{\Cat}D
\]
is a nerve weak equivalence for every functor $D:A\to\Cat$
if and only if every connected component of $A$ is filtered.
The empty index is included. In particular the sequential indices
in the cited manuscript passages satisfy the criterion, whereas
the opposite type category does not.
\end{theorem}
\begin{proof}
\textbf{The filtered implication and reduction to one component.}
For filtered $A$, finite strings, their composition equations,
and finite simplicial homotopy data occur at one stage. Thus
nerve commutes with the colimit and the homotopy-colimit comparison
is a weak equivalence, as proved in Proposition~\ref{F:C22-types}.
Thomason's Grothendieck comparison identifies its source with
$B\int_A D$. All the constructions split as disjoint unions over
components, including the empty case. For necessity it is therefore
enough to take $A$ connected and nonempty. The constant one-point
diagram gives $BA\simeq *$.

\textbf{Common successors, checking the set-valued detection.}
For a discrete-valued diagram the conclusion says that each
component of its category of elements is weakly contractible.
If $a,b$ have no common successor, let $U_a,U_b$ be their principal
cosieves. They are disjoint nonempty connected full subcategories:
every object in each is reached from its named object. Form the
two-point diagram which identifies its two points exactly on
$U_a\amalg U_b$. Its element category is two copies of $A$ glued
along those two cosieves. This also describes the nerves as a union
of subcomplexes, because after a nerve string enters a cosieve it
cannot leave. The two copies have contractible realization. The
kernel of
$H_0(BU_a\amalg BU_b)\to H_0(BA)\oplus H_0(BA)$ is a copy of
$\mathbb Z$, so Mayer--Vietoris gives nonzero first homology in
the connected union. Its ordinary colimit is a singleton, a
contradiction. Every pair of objects has a common successor.

\textbf{The group-valued diagram for a nonequalizable pair.}
Suppose $u,v:p\rightrightarrows q$ have no post-equalizer. Write
$P=A(p,-)$ and $Q=A(q,-)$. At $x$, make a bipartite graph whose
two vertex colors are $Q(x)$ and $P(x)$ and whose edges join
$f\in Q(x)$ to $fu$ and to $fv$. These two edges are distinct:
an equality would itself be a post-equalizer. Let $G(x)$ have
the group presentation
\[
 \left\langle Q(x)\amalg P(x)\ \middle|\
       [f,fu]=[f,fv]=1\text{ for }f\in Q(x)\right\rangle.
\]
Postcomposition preserves the colors and these relations, so this
is an actual group-valued functor on $A$.

\textbf{The primary theorem and every required hypothesis.}
For a finite graph the Salvetti complex is aspherical by
\cite[Theorem~2.6 and Section~2.7]{charney2007}; its signed vertex
links are flag, and its universal cover is contractible. Here the
graph is bipartite, so its only positive-dimensional cells are
the generator circles and one torus cell for each graph edge.
Their cellular boundaries are zero. The same conclusions for an
infinite graph follow by the union of its finite full subgraphs:
each sphere and nullhomotopy uses finitely many cells. Orient every
torus by the $Q$ coordinate followed by the $P$ coordinate.
Transitions preserve these colors, so they send each edge class
to its image edge with coefficient $+1$, even when distinct
vertices of the same color are identified. Therefore, as actual
homology diagrams,
\[
 H_0(BG(-))=\mathbb Z,\qquad
 H_1(BG(-))=\mathbb Z[P\amalg Q],\qquad
 H_2(BG(-))=\mathbb Z[Q\amalg Q],
\]
with all higher groups zero. This checks the naturality needed
for the argument; a list of objectwise homology groups alone
would not suffice.

\textbf{The two specified second homology groups disagree.}
Free abelian representables are projective in the category of
covariant abelian-group diagrams and have colimit $\mathbb Z$.
The derived colimits of the constant diagram are the homology of
$BA$, which is contractible. Thus in the bar double-complex
spectral sequence for $\hocolim_A BG(x)$ every positive
derived-colimit column vanishes. The displayed degree-two row
gives
\[
 H_2\bigl(\hocolim_A BG(x);\mathbb Z\bigr)=\mathbb Z^2.
\]
The ordinary categorical colimit of the one-object groups is
their one-object colimit group: connectedness identifies the
objects and each inverse equation survives the colimit
presentation. Both representable sets have singleton colimit,
so that group has exactly two generators and their commutator
relation. It is $\mathbb Z^2$, whose classifying space has second
homology $\mathbb Z$. The canonical comparison cannot be a weak
equivalence. This contradiction proves post-equalization of every
parallel pair, and hence filteredness.

\textbf{The inventory in the manuscript.}
The sequences have the actual index $\mathbb N$; using a fibred
description of a total category does not reverse that colimit
index. The opposite type category fails filteredness by the two
outer-position maps checked in Proposition~\ref{F:C22-types}.
The stated interval pushout has the direct free-category
universal property verified there, and its two intervals, common
endpoint and concatenation are contractible. This proves that
particular homotopy-pushout comparison without asserting the
universal property for every diagram on a span. These observations
check every use of C.27 in C.28.
\end{proof}

\begin{proposition}[An independent groupoid-suspension replacement for C.15]
\label{F:C15-groupoid-suspension}
Let $G$ be a nontrivial group.  There is a marked ordinary category
with a globally invariant diagram-localization presheaf on all small
tests whose infinity localization has a mapping space $BG$.
\end{proposition}
\begin{proof}
Let $\mathcal M$ have objects $a,x,b$, with $\operatorname{End}(x)=G$,
identity endomorphisms at $a,b$, and unique further arrows
$u:x\to a$ and $v:x\to b$.  Thus $ug=u$ and $vg=v$.
Mark $u$, all elements of $G$, and the identities, but not $v$.

Fix a path type $\tau$ with vertices $0,\ldots,L$.  Its negative
horizontal arrows must be marked; its interior vertical arrows are
arbitrary, and the endpoint vertical arrows are identities.
Write $K_\tau(a,b)$ for this path category.  In any such path an
entry $a\to x$ uses a negative $u$, an exit $x\to a$ uses a
positive $u$, and an exit $x\to b$ uses a positive $v$.
Once $b$ occurs it persists, since a negative $v$ is forbidden.

First normalize the horizontal group labels.  On each maximal
consecutive $x$ interval, set the gauge at its first vertex to $1$
and transport along successive horizontal arrows, using inverses
where their orientations require them.  Conjugating vertical
components by these gauges defines a normalization functor and a
natural isomorphism from the identity to it.  Its image is the full
category of paths whose horizontal group labels are all $1$.

For a normalized path $p$, let $I(p)$ be its last nonempty interval
of $x$ vertices.  Replace every earlier $x$ interval by $a$ and
denote the resulting path by $T(p)$.  For any ladder $p\to q$ one has
\[
 I(q)\subseteq I(p).
\]
Indeed source $a$ vertices must remain $a$, and source $b$ vertices
must remain $b$.  The last target $x$ interval lies after the last
source $a$ vertex: otherwise the target would already have become
$b$ there.  It also lies before the first source $b$ vertex.  The
source vertices strictly between these two vertices are exactly
$I(p)$.  On the retained target interval the normalized vertical
$G$ components are equal.  Use that component for $T(p)\to T(q)$;
all other components are forced identities, $u$'s or $v$'s.
Containment proves compatibility with composition at a common
vertex.  The other naturality equations are $ug=u$ and $vg=v$.
Thus $T$ is a functor, and the components $u$ on the discarded
intervals give a natural transformation $\mathrm{id}\Rightarrow T$.

The image of $T$ is the category $\mathbf BG\times P_\tau$, where
$P_\tau$ consists of admissible nonempty $x$ intervals, ordered by
reverse inclusion.  Admissibility means that the entry edge is
negative and the exit edge is positive.  A ladder between two such
paths is exactly one freely chosen element of $G$ on the smaller
interval, which proves the product description including composition.
When the stage is nonempty, $P_\tau$ has initial object
\[
 I_{\max}=[\text{first negative edge},
                 \text{last positive edge}-1].
\]
It contains every admissible interval.  Projection to $\mathbf BG$
and inclusion at this initial object therefore give a specified
categorical homotopy equivalence.  Combined with the two preceding
normalizations, this proves
\[
 K_\tau(a,b)\simeq\mathbf BG\quad\text{in }\mathrm{Cat}_{\mathrm{hot}}.
\]
These finite zigzags of natural transformations remain such after
$\operatorname{Fun}(Y,-)$ for every small $Y$.

The specified comparisons also control transitions.  On the canonical
$\mathbf BG$ inclusion all horizontal group labels are identities.
Path pullbacks, subdivisions, identity insertions and permitted
compositions retain those identities and retain a vertical element
$g$ at every surviving $x$ vertex.  Subsequent normalization and
last-interval contraction therefore return that same $g$.  All
actual transitions induce the identity of $\mathbf BG$ in
$\mathrm{Cat}_{\mathrm{hot}}$.  The nonempty stages are cofinal.
Using the established ordinary diagram-localization path formula,
\[
 H_{a,b}(Y)=\colim_\tau\pi_0\operatorname{Fun}(Y,K_\tau(a,b))
          =\pi_0\operatorname{Fun}(Y,\mathbf BG)
\]
for every small $Y$.  Functors to a groupoid factor through
$\Pi_1Y$, so this functor inverts every nerve weak equivalence,
without restricting tests to finite categories.

This verifies global invariance for all endpoint pairs as well.
Every $a$-to-$a$ stage has the constant $a$ path as terminal object;
$b$ cannot occur and the needed components are the unique $u$'s.
Every $b$-to-$b$ stage is a point.  There are no $b$-to-$a$ paths,
giving the empty value for nonempty $Y$ and a singleton for empty
$Y$.  Finally the marked natural transformation $u_Y:x_Y\to a_Y$
is inverted in each diagram localization, identifying every
remaining endpoint pair naturally with one of these pairs.

For the infinity-localization computation no general theorem about
arbitrary vertical ladders is needed.  The nerve $N\mathcal M$ is
the union of two right cones on $N\mathbf BG$, with cone vertices
$a,b$ and intersection $N\mathbf BG$.  Both inclusions are
monomorphisms, so this cofibrant span is a homotopy pushout in the
Joyal model.  A functor to an infinity category $\mathcal D$ is
therefore a functor $F:BG\to\mathcal D$, objects $A,B$, and cones
$\alpha:F\to\operatorname{const}A$ and
$\beta:F\to\operatorname{const}B$.  Inverting $u$ means that
$\alpha$ is a pointwise equivalence.  For fixed $A$ the choices
of $F$ with this specified equivalence are contractible.  The
remaining datum has space
\[
 \operatorname{Map}\bigl(BG,\operatorname{Map}_{\mathcal D}(A,B)\bigr).
\]
This is precisely the universal property of the two-object infinity
category with identity endomorphisms, no maps $b\to a$, and
$\operatorname{Map}(a,b)=BG$.  Its mapping component has fundamental
group $G\ne1$ and hence is not discrete.  This proves the desired
example independently of the unavailable original C.15 wording.
For the finite instance $G=C_2$, the program
\path{code/f_groupoid_suspension_table.py} supplies all six arrows,
all defined products, every identity and associativity check, and
closure of the five-arrow marking. The general stage argument above
is independent of any finite bound on the test category.
An independent stronger certificate,
\path{code/b_F_C15_groupoid_suspension.py}, constructs all sixteen
nonempty $a$-to-$b$ stages of length at most four, their actual
normalization and last-interval functors, and the specified natural
transformations. It checks all 25\,408 typed associativity triples,
746 natural-transformation equations, and 52 identity-insertion or
same-sign-composition transition functors. Every induced map on
$C_2$ is the identity. Complete tables and comparisons are stored in
\path{code/outputs/b_F_C15_certificate.json.gz}; the general proof,
rather than this finite bound, establishes the all-small-test claim.
\end{proof}

